\documentclass[titlepage,11pt,a4paper]{book}
\usepackage[a4paper, top=2.5cm, bottom=2.5cm, left=2.5cm, right=2.5cm, headheight=15pt]{geometry}
\usepackage{hyphenat}

% ======================================================================= %
%                      1. PAQUETES BÁSICOS E IDIOMA                       %
% ======================================================================= %
\usepackage[spanish, es-tabla, es-noquoting, es-noindentfirst]{babel}
\addto\captionsspanish{\renewcommand{\tablename}{Tabla}}
\addto\captionsspanish{\renewcommand{\listtablename}{Índice de tablas}}

\usepackage{csquotes}
\usepackage[utf8]{inputenc}
\usepackage[T1]{fontenc}
\usepackage{newpxtext}
\usepackage{newpxmath}        
\usepackage{microtype}      
\linespread{1.05}           
\usepackage{textcomp}     

% --- Nomenclatura (Índice de Símbolos) ---
\usepackage[intoc, spanish]{nomencl}
\makenomenclature
\usepackage{etoolbox}
\renewcommand\nomname{Glosario de Notación}
\renewcommand{\nomlabel}[1]{\textbf{#1}\dotfill}

% --- Índice alfabñetico ---
\usepackage{imakeidx}
\makeindex[intoc, title=Índice Alfabético]

% ======================================================================= %
%                      2. MATEMÁTICAS Y SÍMBOLOS                          %
% ======================================================================= %
\let\Bbbk\relax
\let\openbox\relax 
\usepackage{amsmath,amsfonts,amssymb,amsthm} 
\usepackage{mathtools}
\usepackage{cancel}
\usepackage{commath}
\usepackage{listings}
\usepackage{enumerate}
\usepackage{bussproofs} 


% ======================================================================= %
%                  3. ELEMENTOS DE ESCRITURA Y REFERENCIA                 %
% ======================================================================= %
\renewcommand{\contentsname}{Índice General}
\usepackage{chngcntr}
% --- Control de Viudas y Huérfanas ---
\clubpenalty=10000
\widowpenalty=10000
\displaywidowpenalty=10000
% ======================================================================= %
%                 CONFIGURACIÓN DE ENCABEZADOS Y PIE DE PÁGINA            %
% ======================================================================= %
\usepackage{fancyhdr} 
\pagestyle{fancy}
\fancyhf{} 
\fancyhfoffset[LE,RO]{0pt}

\renewcommand{\headrulewidth}{1pt}
\renewcommand{\footrulewidth}{1pt}

\let\oldheadrule\headrule
\let\oldfootrule\footrule
\renewcommand{\headrule}{{\color{colorAsignatura!60}\oldheadrule}}
\renewcommand{\footrule}{{\color{colorAsignatura!60}\oldfootrule}}

\fancyhead[L]{\textcolor{colorAsignatura}{\small\bfseries\leftmark}}
\fancyhead[R]{\textcolor{colorAsignatura}{\small\scshape\materia}}

\fancyfoot[LE,RO]{%
    \begin{tikzpicture}[baseline=0pt]
        \fill[colorAsignatura] (0,0) rectangle (1, 0.7);
        \node[white, font=\bfseries] at (0.5, 0.35) {\thepage};
    \end{tikzpicture}%
}

\fancypagestyle{plain}{
    \fancyhf{} 
    \renewcommand{\headrulewidth}{0pt} 
    \renewcommand{\footrulewidth}{1pt} 
    \fancyfoot[LE,RO]{%
        \begin{tikzpicture}[baseline=0pt]
            \fill[colorAsignatura] (0,0) rectangle (1, 0.7);
            \node[white, font=\bfseries] at (0.5, 0.35) {\thepage};
        \end{tikzpicture}%
    }
}
\usepackage[normalem]{ulem} 
\usepackage{emptypage} 

% ======================================================================= %
%                      5. TABLAS, FIGURAS Y AÑADIDOS                      %
% ======================================================================= %
\usepackage[table,dvipsnames]{xcolor}
\usepackage{graphicx}
\usepackage{epstopdf}
\usepackage{float}
\usepackage{caption}
\captionsetup{
    labelfont={color=colorAsignatura, bf},
    textfont={small, it},
    margin=10pt
}
\usepackage{subcaption}
\usepackage{booktabs}
\usepackage{tabularx}
\usepackage{makecell}
\usepackage{multirow}
\usepackage{nicematrix}
\usepackage{siunitx}
\sisetup{output-decimal-marker = {,}}
\usepackage{rotating}
\setlength{\floatsep}{15pt} 
\setlength{\textfloatsep}{20pt} 

% --- Espaciado y sangría de párrafos ---
\setlength{\parskip}{0.8em plus 0.2em minus 0.2em} % El salto vertical entre párrafos
\setlength{\parindent}{0pt} % 0pt quita la sangría inicial. Pon 1.5em si prefieres mantenerla.

% ==========================================
% CONFIGURACIÓN DE LISTAS (enumitem)
% ==========================================
\usepackage{enumitem}

% Ajuste global de todas las listas (itemize y enumerate)
\setlist{
  topsep=0.6em,      % Espacio antes de empezar la lista
  itemsep=0.5em,     % Espacio entre un ítem y el siguiente
  parsep=0.2em,      % Espacio entre párrafos dentro de un mismo ítem
  partopsep=0.2em    % Espacio extra si la lista viene precedida de una línea en blanco
}

\setlist[enumerate,1]{label=\textit{\roman*.}} % Mantiene tu estilo de números romanos


%EPÍGRAFES - %
\usepackage{epigraph}
\setlength{\epigraphwidth}{0.6\textwidth}
\setlength{\epigraphrule}{0pt}
\renewcommand{\textflush}{flushright}
% ======================================================================= %
%                      6. ELEMENTOS DE PORTADA                            %
% ======================================================================= %
\newcommand{\titulo}{Recolección de apuntes de} 
\newcommand{\facultad}{Ciencias}
\newcommand{\autor}{Iván Jesús Navarro Ugidos}
\newcommand{\materia}{Análisis Funcional} 
\newcommand{\ciudad}{Oviedo}
\newcommand{\fecha}{\today}

% ======================================================================= %
%                      7. ELEMENTOS DE DIBUJADO (TIKZ)                    %
% ======================================================================= %
\usepackage{tikz}
\usepackage{tikz-cd}
\usepackage{pgfplots}
\usetikzlibrary{babel}
\pgfplotsset{compat=1.18}
\usetikzlibrary{positioning, arrows.meta}
\usetikzlibrary{shapes, decorations, calc, arrows}
\usetikzlibrary{3d,fit,backgrounds, decorations.text}
\usetikzlibrary{positioning, shapes.symbols}
\usetikzlibrary{decorations.pathreplacing, calligraphy}

% 1. Para poder insertar tu portada compilada previamente
\usepackage{pdfpages} 

% 2. Para activar la externalización de gráficos
\usetikzlibrary{external}
\tikzexternalize[prefix=figuras-tikz-func/]
\tikzset{external/only named=true}

% Añade esta línea para proteger a las cajas de tcolorbox


\tikzset{%
    myarrow/.style = {-Stealth, shorten >=5pt}
} 
\newcommand{\mypoint}[2]{\tikz[remember picture]{\node[inner sep=0, anchor=base](#1){$#2$};}}

\tikzset{>=latex}
\usepackage{adjustbox}

%=========ESTILO PARA MUNDOS POSIBLES==========%
\tikzset{
    % Estilo para mundos posibles (nodos)
    mundo/.style={
        circle, 
        draw=colorAsignatura, 
        fill=colorAsignatura!10, 
        line width=1.5pt, 
        minimum size=1cm, 
        font=\bfseries\color{colorAsignatura}
    },
    %=========ESTILO PARA RELACIONES DE ACCESIBILIDAD==========%
    acc/.style={
        -{Stealth[length=3mm, width=2mm]}, 
        colorAsignatura, 
        line width=1.2pt
    }
}
% ======================================================================= %
%                      8. ATAJOS Y COMANDOS PERSONALIZADOS            
% ======================================================================= %
\newcommand{\NN}{\ensuremath{\textcolor{OliveGreen}{\mathbb{N}}}}
\newcommand{\ZZ}{\ensuremath{\textcolor{violet}{\mathbb{Z}}}}
\newcommand{\QQ}{\ensuremath{\textcolor{BurntOrange}{\mathbb{Q}}}}
\newcommand{\II}{\ensuremath{\textcolor{Aquamarine}{\mathbb{I}}}}
\newcommand{\RR}{\ensuremath{\textcolor{blue}{\mathbb{R}}}}
\newcommand{\CC}{\ensuremath{\textcolor{Red}{\mathbb{C}}}}
\newcommand{\filt}{\ensuremath{\textcolor{Orange}{\mathcal{F}}}}
\newcommand{\gilt}{\ensuremath{\textcolor{Magenta}{\mathcal{G}}}}
\newcommand{\ultf}{\ensuremath{\textcolor{Plum}{\mathfrak{U}}}}
\newcommand{\ii}{\ensuremath{\textcolor{TealBlue}{i}}}
\newcommand{\jj}{\ensuremath{\textcolor{OliveGreen}{j}}}
\newcommand{\nn}{\ensuremath{\textcolor{Plum}{n}}}
\newcommand{\mm}{\ensuremath{\textcolor{Magenta}{m}}}
\newcommand{\KK}{\ensuremath{\mathbb{K}}}
\newcommand{\VV}{\ensuremath{\mathcal{V}}}
\newcommand{\alpa}{\textcolor{orange}{\alpha}}
\newcommand{\berta}{\textcolor{blue}{\beta}}
\newcommand{\eps}{\textcolor{ForestGreen}{\epsilon}}
\newcommand{\mum}{\textcolor{Plum}{\mu}}
\newcommand{\delt}{\textcolor{blue}{\delta}}
\newcommand{\Delt}{\textcolor{ForestGreen}{\Delta}}
\newcommand{\etat}{\textcolor{Rhodamine}{\eta}}
\newcommand{\gam}{\textcolor{ForestGreen}{\gamma}}
\newcommand{\lamb}{\textcolor{TealBlue}{\lambda}}
\newcommand{\taut}{\textcolor{Fuchsia}{\tau}}
\newcommand{\kk}{\textcolor{orange}{k}}
\newcommand{\lal}{\textcolor{Purple}{l}}
\newcommand{\tet}{\textcolor{BurntOrange}{t}}
\newcommand*{\temp}{\multicolumn{1}{c|}{0}}


\newcommand{\Chapter}[1]{\chapter{#1}
\begin{mdframed}{\Large\bfseries Índice}
\minitoc
\vspace{-8mm}
\end{mdframed}
}


\definecolor{colorAsignatura}{RGB}{77, 138, 84} 

\usepackage{titlesec}




% --- MACROS AUXILIARES RECUPERADAS ---

\newcommand{\centrarTextoSeccion}[1]{%
    \raisebox{0.05cm}{%
        \parbox[c][2.1cm][c]{\dimexpr\textwidth-1.5cm\relax}{\centering\color{colorAsignatura}#1}%
    }%
}

\newcommand{\centrarTextoCapitulo}[1]{%
    \raisebox{-1cm}{%
        \parbox[c][4cm][c]{\dimexpr\textwidth-1.5cm\relax}{\centering\color{colorAsignatura}#1}%
    }%
}

% ==========================================
% CONFIGURACIÓN DE SECCIONES
% ==========================================

\DeclareRobustCommand{\dibujarSeccion}[1]{%
    \begin{tikzpicture}[remember picture, external/export = false]
        % 1. El texto se alinea exactamente con el margen izquierdo normal (x=0). 
        %    Ocupa todo el ancho hasta el margen derecho (\textwidth).
        \node[text width=\textwidth, align=left, inner ysep=12pt, inner xsep=0pt] 
            (sectext) at (0, 0) [anchor=west] {#1};
            
        \begin{scope}[overlay]
            % 2. La caja azul nace en el extremo físico izquierdo del papel y se CORTA 
            %    exactamente 0.4 cm antes de chocar con el texto.
            \fill[colorAsignatura] 
                (current page.west |- sectext.south) rectangle ([xshift=-0.4cm]sectext.north west);
            
            % 3. El número se desplaza a la izquierda, anclado dentro de la franja azul,
            %    dejando un pequeño respiro respecto al corte.
            \node[white, anchor=east] at ([xshift=-0.5cm]sectext.west) 
                {\fontsize{34}{34}\selectfont\thesection};
            
            % 4. Las líneas horizontales nacen en el corte de la franja azul y mueren
            %    exactamente en el margen derecho del texto.
            \draw[colorAsignatura, line width=2pt] 
                ([xshift=-0.4cm, yshift=-1pt]sectext.north west) -- (sectext.north east);
            \draw[colorAsignatura, line width=2pt] 
                ([xshift=-0.4cm, yshift=1pt]sectext.south west) -- (sectext.south east);
        \end{scope}
    \end{tikzpicture}%
}



\titleformat{\section}[block]
{\normalfont\huge\bfseries\color{colorAsignatura}} 
{}
{0pt}
{\thispagestyle{plain}\dibujarSeccion}

% Compensamos el espacio vertical directamente aquí. 
% \titlespacing*{comando}{izq}{espacio ANTES}{espacio DESPUÉS}
\titlespacing*{\section}{0pt}{6ex}{4ex}


\DeclareRobustCommand{\dibujarSeccionEstrella}[1]{%
    \begin{tikzpicture}[remember picture, external/export = false]
        \node[text width=\textwidth, align=left, inner ysep=12pt, inner xsep=0pt] 
            (sectext) at (0, 0) [anchor=west] {#1};
            
        \begin{scope}[overlay]
            \fill[colorAsignatura] 
                (current page.west |- sectext.south) rectangle ([xshift=-0.4cm]sectext.north west);
            
            % HEMOS ELIMINADO EL NODO DE \thesection AQUÍ
            
            \draw[colorAsignatura, line width=2pt] 
                ([xshift=-0.4cm, yshift=-1pt]sectext.north west) -- (sectext.north east);
            \draw[colorAsignatura, line width=2pt] 
                ([xshift=-0.4cm, yshift=1pt]sectext.south west) -- (sectext.south east);
        \end{scope}
    \end{tikzpicture}%
}

\titleformat{name=\section,numberless}[block]
{\normalfont\huge\bfseries\color{colorAsignatura}} 
{}
{0pt}
{\thispagestyle{plain}\dibujarSeccionEstrella}

\titlespacing*{\section}{0pt}{6ex}{4ex}
% ==========================================
% CONFIGURACIÓN DE SUBSECCIONES (Limpieza editorial)
% ==========================================
\titleformat{\subsection}
  {\normalfont\Large\bfseries\color{colorAsignatura!85!black}} % Ligeramente más oscuro que la sección para jerarquía
  {\thesubsection} % El número
  {1em} % Espacio entre número y texto
  {}

% Control de espaciados: \titlespacing*{comando}{izq}{arriba}{abajo}
\titlespacing*{\subsection}{0pt}{4.5ex plus 1ex minus .2ex}{2ex plus .2ex}

% ==========================================
% CONFIGURACIÓN DE CAPÍTULOS
% ==========================================

% --- Capítulos Numerados ---
\titleformat{\chapter}[hang]
{\normalfont\huge\bfseries} 
{%
    \hbox to 1.5cm{%
        \begin{tikzpicture}[remember picture, overlay,  external/export = false]
            \fill[colorAsignatura] (current page.north west) rectangle ++(3.5cm,-6cm);
            \node[white] at ([xshift=1.75cm, yshift=-2.2cm]current page.north west) {\Large\scshape Capítulo};
            \node[white] at ([xshift=1.75cm, yshift=-4.2cm]current page.north west) {\fontsize{75}{75}\selectfont\thechapter};
            \draw[colorAsignatura, line width=3pt, yshift=1-5pt] ([xshift=3.5cm, yshift=-6cm+1.5pt]current page.north west) -- ([yshift=-6cm]current page.north east);
        \end{tikzpicture}\hss
    }%
}
{0pt}
{\fontsize{40}{40}\selectfont\centrarTextoCapitulo}

% --- NUEVA Macro para Capítulos NO Numerados ---
\DeclareRobustCommand{\dibujarTituloCapituloEstrella}[1]{%
    \begin{tikzpicture}[remember picture, overlay,  external/export = false]
        \fill[colorAsignatura] (current page.north west) rectangle ([yshift=-3.5cm]current page.north east);
        \node[white, align=center, text width=\paperwidth] 
            at ([yshift=-1.75cm]current page.north) {
                \fontsize{40}{45}\selectfont\bfseries #1
            };
    \end{tikzpicture}%
}


% --- Capítulos NO Numerados (Índice, Bibliografía...) ---
\titleformat{name=\chapter,numberless}[display]
{\normalfont\huge\bfseries} 
{}
{0pt}
{\dibujarTituloCapituloEstrella}

% --- Espaciado General ---
\titlespacing*{\chapter}{0pt}{-1.5cm}{3.2cm}


% ==========================================
% CONFIGURACIÓN DE PARTES (Interlineado y sin guiones)
% ==========================================
\newcommand{\dibujarTituloParte}[1]{%
    \begin{tikzpicture}[remember picture, overlay,  external/export = false]
        % Fondo lateral
        \fill[colorAsignatura] (current page.north west) rectangle ([xshift=4cm]current page.south west);
        \node[white, rotate=90, anchor=center] at ([xshift=2cm]current page.west) {
            \fontsize{60}{60}\selectfont\bfseries PARTE \Roman{part}
        };
        
        % NODO CENTRAL: 
        % Al poner \fontsize en las opciones del nodo (font=...), TikZ asume 
        % el control total y aplica los 60pt de separación entre las líneas sin fallar.
        \node[text width=\dimexpr\paperwidth-6cm\relax, 
              align=center, 
              inner ysep=1.5cm,
              font=\fontsize{40}{60}\selectfont\bfseries\color{colorAsignatura},
              execute at begin node=\setlength{\baselineskip}{60pt}] 
            (textoparte) at ([xshift=2cm]current page.center) {
                % Usamos tu paquete 'hyphenat' para blindar el texto contra los guiones
                \nohyphens{\MakeUppercase{#1}}
            };
            
        % Líneas adaptativas superior e inferior
        \draw[colorAsignatura, line width=3pt] 
            ([xshift=4cm]textoparte.north west -| current page.west) -- ([xshift=-1cm]textoparte.north east -| current page.east);
        \draw[colorAsignatura, line width=3pt] 
            ([xshift=4cm]textoparte.south west -| current page.west) -- ([xshift=-1cm]textoparte.south east -| current page.east);
    \end{tikzpicture}%
}

\titleformat{\part}[display]
{\thispagestyle{empty}} 
{} 
{0pt}
{\dibujarTituloParte} 

\titlespacing*{\part}{0pt}{0pt}{0pt}

% ======================================================================= %
%                      9. REFERENCIAS INTERNAS Y EXTERNAS         
% ======================================================================= %
\usepackage{varioref}
\usepackage[colorlinks=true, 
            linkcolor=colorAsignatura!80!black, 
            urlcolor=colorAsignatura!80!black, 
            citecolor=colorAsignatura!80!black]{hyperref}
\usepackage[nameinlink, spanish]{cleveref}
\crefname{thm}{Teorema}{Teoremas}
\crefname{defn}{Definición}{Definiciones}
\crefname{lem}{Lema}{Lemas}
\crefname{cor}{Corolario}{Corolarios}
\crefname{prop}{Proposición}{Proposiciones}
\crefname{axm}{Axioma}{Axiomas}
\crefname{rem}{Nota}{Notas}

\usepackage{minitoc}
\mtcsetdepth{minitoc}{1}
\mtcsettitle{minitoc}{Índice del Capítulo}
\mtcsetrules{minitoc}{on}
\nomtcpagenumbers

% --- Configuración Minitoc ---
\setcounter{minitocdepth}{1} 
\mtcsetrules{minitoc}{off} 
\mtcsettitle{minitoc}{Índice del Capítulo}
\mtcsetfont{minitoc}{*}{\large\scshape\color{colorAsignatura}} 
\mtcsetfont{minitoc}{section}{\normalfont\small\color{colorAsignatura!80!black}}
\usepackage{mdframed}
\renewcommand{\mtctitle}{Índice}
\usepackage[backend=biber, style=alphabetic, sorting=nyt]{biblatex}

% Cambiamos el título por defecto a español
\DefineBibliographyStrings{spanish}{
  bibliography = {Bibliografía},
}

% Enlazamos el archivo externo donde guardaremos los libros
\addbibresource{bibliografia.bib}

\usepackage{bookmark}
\bookmarksetup{
  numbered,
  open=true,
  openlevel=1
}

% ======================================================================= %
%                           10. TEOREMAS Y CAJAS                           %
% ======================================================================= %
\usepackage[most]{tcolorbox}
\tcbuselibrary{skins}
\tcbuselibrary{theorems}
\tcbuselibrary{listings}
\tcbset{shield externalize}
\swapnumbers

% Comando para rodear un símbolo con borde negro grueso y espacio a la derecha
\newcommand{\borde}[1]{%
  \begingroup
  \setlength{\fboxsep}{2.5pt}% Separación interna entre el símbolo y el borde
  \setlength{\fboxrule}{1pt}% Grosor del borde negro
  \textcolor{black}{\fbox{\normalcolor\ensuremath{#1}}}% Dibuja la caja
  \hspace{4pt}% <--- LIGERO ESPACIO HORIZONTAL A LA DERECHA (puedes variar los pt)
  \vspace{0.5 pt}
  \endgroup
}


\newtcbtheorem[number within=chapter]{defn}{Definición}%
{label type=defn,colback=colorAsignatura!10,title=Definition,
	attach boxed title to top left={xshift=-4mm},boxrule=0pt,after skip=1cm,before skip=1cm,right skip=0cm,breakable,fonttitle=\bfseries,toprule=0pt,bottomrule=0pt,rightrule=0pt,leftrule=4pt,arc=0mm,skin=enhancedlast jigsaw,rounded corners,colframe=colorAsignatura,colbacktitle=colorAsignatura!80,boxed title style={
		frame code={ 
			\fill[colorAsignatura!80] (frame.south west)-- (frame.north west)--(frame.north east)to [in=-1,out=1,looseness=2] (frame.south east)--cycle;
            \draw[line width=0.8mm,colorAsignatura!80]([xshift=3mm]frame.north east)to [in=-1,out=1,looseness=2]([xshift=3mm]frame.south east);
			\draw[line width=0.8mm,colorAsignatura!80]([xshift=7mm]frame.north east)to [in=-1,out=1,looseness=2]([xshift=7mm]frame.south east);
            \draw[line width=0.8mm,colorAsignatura!80]([xshift=11mm]frame.north east)to [in=-1,out=1,looseness=2]([xshift=11mm]frame.south east);
			\fill[colorAsignatura!40!black](frame.south west)--+(4mm,-2mm)--+(4mm,2mm)--cycle;
		}
	}
}{defn}
\NewTcbTheorem[number within=chapter]{thm}{Teorema}
{label type=thm,
    colback=colorAsignatura!10, enhanced, title=Teorema,attach boxed title to top left={xshift=-4mm}, boxrule=0pt,
    after skip=1cm, before skip=1cm, right skip=0cm, breakable, fonttitle=\bfseries, toprule=0pt, bottomrule=0pt, rightrule=0pt, leftrule=4pt, arc=0mm, skin=enhancedlast jigsaw, sharp corners, colframe=colorAsignatura, colbacktitle=colorAsignatura,
    boxed title style={
        frame code={ 
            \fill[colorAsignatura]
                (frame.south west)--(frame.north west)--(frame.north east)--
                ([xshift=3mm]frame.east)--(frame.south east)--cycle;
            \draw[line width=1mm, colorAsignatura]
                ([xshift=2mm]frame.north east)--([xshift=5mm]frame.east)--([xshift=2mm]frame.south east);
            \draw[line width=1mm, colorAsignatura]
                ([xshift=5mm]frame.north east)--([xshift=8mm]frame.east)--([xshift=5mm]frame.south east);
            \draw[line width=1mm, colorAsignatura]
                ([xshift=8mm]frame.north east)--([xshift=11mm]frame.east)--([xshift=8mm]frame.south east);
            \fill[colorAsignatura!40!black]
                (frame.south west)--+(4mm,-2mm)--+(4mm,2mm)--cycle;
        }
    }
}{thm}

\NewTcbTheorem[number within=chapter]{thmroj}{Teorema}
{label type=thm,colback=red!20!white,enhanced,title=Teorema,
	attach boxed title to top left={xshift=-4mm},boxrule=0pt,after skip=1cm,before skip=1cm,right skip=0cm,breakable,fonttitle=\bfseries,	toprule=0pt,bottomrule=0pt,rightrule=0pt,leftrule=4pt,arc=0mm,skin=enhancedlast jigsaw,sharp corners,colframe=red!60!black,colbacktitle=red!80!black,boxed title style={
		frame code={ 
			\fill[red!80!black](frame.south west)--(frame.north west)--(frame.north east)--([xshift=3mm]frame.east)--(frame.south east)--cycle;
			\draw[line width=1mm,red!80!black]([xshift=2mm]frame.north east)--([xshift=5mm]frame.east)--([xshift=2mm]frame.south east);
			\draw[line width=1mm,red!80!black]([xshift=5mm]frame.north east)--([xshift=8mm]frame.east)--([xshift=5mm]frame.south east);
            \draw[line width=1mm,red!80!black]([xshift=8mm]frame.north east)--([xshift=11mm]frame.east)--([xshift=8mm]frame.south east);
			\fill[red!60!white](frame.south west)--+(4mm,-2mm)--+(4mm,2mm)--cycle;
		}
	}
}{teoroj}

\NewTcbTheorem[number within=chapter]{prop}{Proposición}
{label type=prop,colback=colorAsignatura!10, enhanced, title=Teorema,attach boxed title to top left={xshift=-4mm}, boxrule=0pt,
    after skip=1cm, before skip=1cm, right skip=0cm, breakable, fonttitle=\bfseries, toprule=0pt, bottomrule=0pt, rightrule=0pt, leftrule=4pt, arc=0mm, skin=enhancedlast jigsaw, sharp corners, colframe=colorAsignatura, colbacktitle=colorAsignatura,
    boxed title style={
        frame code={ 
            \fill[colorAsignatura]
                (frame.south west)--(frame.north west)--(frame.north east)--
                ([xshift=3mm]frame.east)--(frame.south east)--cycle;
            \draw[line width=1mm, colorAsignatura]
                ([xshift=2mm]frame.north east)--([xshift=5mm]frame.east)--([xshift=2mm]frame.south east);
            \draw[line width=1mm, colorAsignatura]
                ([xshift=5mm]frame.north east)--([xshift=8mm]frame.east)--([xshift=5mm]frame.south east);
            \draw[line width=1mm, colorAsignatura]
                ([xshift=8mm]frame.north east)--([xshift=11mm]frame.east)--([xshift=8mm]frame.south east);
            \fill[colorAsignatura!40!black]
                (frame.south west)--+(4mm,-2mm)--+(4mm,2mm)--cycle;
        }
    }
}{prop}

\NewTcbTheorem[number within=chapter]{proproj}{Proposición}
{label type=prop,colback=orange!20!white,enhanced,title=Teorema,
	attach boxed title to top left={xshift=-4mm},boxrule=0pt,after skip=1cm,before skip=1cm,right skip=0cm,breakable,fonttitle=\bfseries,	toprule=0pt,bottomrule=0pt,rightrule=0pt,leftrule=4pt,arc=0mm,skin=enhancedlast jigsaw,sharp corners,colframe=orange!85!black,colbacktitle=orange!95!black,boxed title style={
		frame code={ 
			\fill[orange!95!black](frame.south west)--(frame.north west)--(frame.north east)--([xshift=3mm]frame.east)--(frame.south east)--cycle;
			\draw[line width=1mm,orange!95!black]([xshift=2mm]frame.north east)--([xshift=5mm]frame.east)--([xshift=2mm]frame.south east);
			\draw[line width=1mm,orange!95!black]([xshift=5mm]frame.north east)--([xshift=8mm]frame.east)--([xshift=5mm]frame.south east);
			\fill[orange!60!white](frame.south west)--+(4mm,-2mm)--+(4mm,2mm)--cycle;
		}
	}
}{proproj}

\NewTcbTheorem[number within=chapter]{prop+}{Propiedades}
{label type=prop,colback=colorAsignatura!10,enhanced,title=Teorema,
	attach boxed title to top left={xshift=-4mm},boxrule=0pt,after skip=1cm,before skip=1cm,right skip=0cm,breakable,fonttitle=\bfseries,	toprule=0pt,bottomrule=0pt,rightrule=0pt,leftrule=4pt,arc=0mm,skin=enhancedlast jigsaw,sharp corners,colframe=colorAsignatura,colbacktitle=colorAsignatura,boxed title style={
		frame code={ 
			\fill[colorAsignatura](frame.south west)--(frame.north west)--(frame.north east)--([xshift=3mm]frame.east)--(frame.south east)--cycle;
			\draw[line width=1mm,colorAsignatura]([xshift=2mm]frame.north east)--([xshift=5mm]frame.east)--([xshift=2mm]frame.south east);
			\fill[colorAsignatura!40!black](frame.south west)--+(4mm,-2mm)--+(4mm,2mm)--cycle;
		}
	}
}{prop+}

\NewTcbTheorem[number within=chapter]{proproj+}{Propiedades}
{label type=prop,colback=Dandelion!20!white,enhanced,title=Teorema,
	attach boxed title to top left={xshift=-4mm},boxrule=0pt,after skip=1cm,before skip=1cm,right skip=0cm,breakable,fonttitle=\bfseries,	toprule=0pt,bottomrule=0pt,rightrule=0pt,leftrule=4pt,arc=0mm,skin=enhancedlast jigsaw,sharp corners,colframe=Dandelion!85!black,colbacktitle=Dandelion!95!black,boxed title style={
		frame code={ 
			\fill[Dandelion!95!black](frame.south west)--(frame.north west)--(frame.north east)--([xshift=3mm]frame.east)--(frame.south east)--cycle;
			\draw[line width=1mm,Dandelion!95!black]([xshift=2mm]frame.north east)--([xshift=5mm]frame.east)--([xshift=2mm]frame.south east);
			\fill[Dandelion!60!white](frame.south west)--+(4mm,-2mm)--+(4mm,2mm)--cycle;
		}
	}
}{proproj+}

\NewTcbTheorem[number within=chapter]{cor}{Corolario}
{label type=cor,colback=colorAsignatura!10,enhanced,title=Teorema,
	attach boxed title to top left={xshift=-4mm},boxrule=0pt,after skip=1cm,before skip=1cm,right skip=0cm,breakable,fonttitle=\bfseries,	toprule=0pt,bottomrule=0pt,rightrule=0pt,leftrule=4pt,arc=0mm,skin=enhancedlast jigsaw,sharp corners,colframe=colorAsignatura,colbacktitle=colorAsignatura,boxed title style={
		frame code={ 
			\fill[colorAsignatura](frame.south west)--(frame.north west)--(frame.north east)--([xshift=3mm]frame.east)--(frame.south east)--cycle;
			\draw[line width=1mm,colorAsignatura]([xshift=2mm]frame.north east)--([xshift=5mm]frame.east)--([xshift=2mm]frame.south east);
			\fill[colorAsignatura!40!black](frame.south west)--+(4mm,-2mm)--+(4mm,2mm)--cycle;
		}
	}
}{cor}

\NewTcbTheorem[number within=chapter]{lem}{Lema}
{label type=lem,colback=colorAsignatura!10,enhanced,title=Teorema,
	attach boxed title to top left={xshift=-4mm},boxrule=0pt,after skip=1cm,before skip=1cm,right skip=0cm,breakable,fonttitle=\bfseries,	toprule=0pt,bottomrule=0pt,rightrule=0pt,leftrule=4pt,arc=0mm,skin=enhancedlast jigsaw,sharp corners,colframe=colorAsignatura,colbacktitle=colorAsignatura,boxed title style={
		frame code={ 
			\fill[colorAsignatura](frame.south west)--(frame.north west)--(frame.north east)--([xshift=3mm]frame.east)--(frame.south east)--cycle;
			\draw[line width=1mm,colorAsignatura]([xshift=2mm]frame.north east)--([xshift=5mm]frame.east)--([xshift=2mm]frame.south east);
			\fill[colorAsignatura!40!black](frame.south west)--+(4mm,-2mm)--+(4mm,2mm)--cycle;
		}
	}
}{lem}

\NewTcbTheorem[number within=chapter]{axm}{Axioma}
{label type=axm,
    colback=red!20!white, 
    enhanced, 
    title=Axioma,
    attach boxed title to top left={xshift=-4mm}, 
    boxrule=0pt,
    after skip=1cm, 
    before skip=1cm, 
    right skip=0cm, 
    breakable, 
    fonttitle=\bfseries, 
    toprule=0pt, 
    bottomrule=0pt, 
    rightrule=0pt, 
    leftrule=4pt, 
    arc=0mm, 
    skin=enhancedlast jigsaw, 
    sharp corners, 
    colframe=red!60!black, 
    colbacktitle=red!80!black,
    boxed title style={
        frame code={ 
            % Cuerpo del título en bloque recto
            \fill[red!80!black]
                (frame.south west)--(frame.north west)--(frame.north east)--(frame.south east)--cycle;
            
            % Doble pilar vertical a la derecha (diseño de axioma)
            \draw[line width=1.2mm, red!80!black]
                ([xshift=2mm]frame.north east)--([xshift=2mm]frame.south east);
            \draw[line width=1.2mm, red!80!black]
                ([xshift=5mm]frame.north east)--([xshift=5mm]frame.south east);
            
            % Triángulo de relieve en paralelo (idéntico en dimensiones y tono al teorema rojo)
            \fill[red!40!black]
                (frame.south west)--+(4mm,-2mm)--+(4mm,2mm)--cycle;
        }
    }
}{ax}

\NewTcbTheorem[number within=chapter]{aux}{Resultado Auxiliar}
{colback=Blue!20!white,enhanced,title=Teorema,
	attach boxed title to top left={xshift=-4mm},boxrule=0pt,after skip=1cm,before skip=1cm,right skip=0cm,breakable,fonttitle=\bfseries,	toprule=0pt,bottomrule=0pt,rightrule=0pt,leftrule=4pt,arc=0mm,skin=enhancedlast jigsaw,sharp corners,colframe=Blue!85!black,colbacktitle=Blue!95!black,boxed title style={
		frame code={ 
			\fill[Blue!95!black](frame.south west)--(frame.north west)--(frame.north east)--([xshift=3mm]frame.east)--(frame.south east)--cycle;
			\draw[line width=1mm,Blue!95!black]([xshift=2mm]frame.north east)--([xshift=5mm]frame.east)--([xshift=2mm]frame.south east);
			\fill[Blue!60!white](frame.south west)--+(4mm,-2mm)--+(4mm,2mm)--cycle;
		}
	}
}{aux}

\newtcolorbox{auxi}{
  colback=Blue!20!white,
  enhanced,
  boxrule=0pt,
  after skip=1cm,
  before skip=1cm,
  right skip=0cm,
  breakable,
  fonttitle=\bfseries,
  toprule=0pt,
  bottomrule=0pt,
  rightrule=0pt,
  leftrule=4pt,
  arc=0mm,
  skin=enhancedlast jigsaw,
  sharp corners,
  colframe=Blue!85!black
}


\newtcbtheorem[number within=chapter]{rem}{Nota}%
{label type=rem,colback=gray!20!white,title=Nota,
	attach boxed title to top left={xshift=-4mm},boxrule=0pt,after skip=1cm,before skip=1cm,right skip=0cm,breakable,fonttitle=\bfseries,toprule=0pt,bottomrule=0pt,rightrule=0pt,leftrule=4pt,arc=0mm,skin=enhancedlast jigsaw,rounded corners,colframe=gray!80!white,colbacktitle=gray!70!white,boxed title style={
		frame code={ 
			\fill[gray!70!white] (frame.south west)-- (frame.north west)--(frame.north east)to [in=-1,out=1,looseness=2] (frame.south east)--cycle;
            \draw[line width=0.8mm,gray!70!white]([xshift=3mm]frame.north east)to [in=-1,out=1,looseness=2]([xshift=3mm]frame.south east);
			\fill[gray!30!white](frame.south west)--+(4mm,-2mm)--+(4mm,2mm)--cycle;
		}
	}
}{rem}

% ==========================================
% ENTORNO DE EJEMPLOS 
% ==========================================
\NewTcbTheorem[number within=chapter]{ejemplo}{Ejemplo}
{label type=ejm,
    colback=white, enhanced, title=Ejemplo,
    attach boxed title to top left={xshift=-4mm}, boxrule=0pt,
    after skip=1cm, before skip=1cm, right skip=0cm, breakable, 
    fonttitle=\bfseries\sffamily, coltitle=black, % Título negro
    toprule=0pt, bottomrule=0pt, rightrule=0pt, leftrule=2pt, arc=0mm, 
    skin=enhancedlast jigsaw, sharp corners, 
    colframe=black!40, colbacktitle=black!10, % Tonos grises
    boxed title style={
        frame code={ 
            % Caja plana y gris, sin los adornos puntiagudos
            \fill[black!10] (frame.south west)--(frame.north west)--(frame.north east)--(frame.south east)--cycle;
            % Pliegue izquierdo inferior para mantener la filosofía global del documento
            \fill[black!50] (frame.south west)--+(4mm,-2mm)--+(4mm,2mm)--cycle;
        }
    }
}{ejm}

% ======================================================================= %
%                              DEMOSTRACIONES                             %
% ======================================================================= %
\newtcolorbox{dem}[1][Demostración]{
  enhanced,
  breakable,
  frame hidden, 
  boxrule=0pt,
  % Línea casi negra con un matiz del color de la asignatura
  borderline west={3.5pt}{0pt}{colorAsignatura!25!black}, 
  % Fondo muy claro, blanco manchado levemente con el color de la asignatura
  colback=colorAsignatura!3!white, 
  coltitle=colorAsignatura!25!black, 
  % --- Cambio tipográfico: Serif clásica, negrita y versalitas ---
  fonttitle=\rmfamily\bfseries\scshape\large, 
  title={#1}, 
  before skip=3mm, 
  after skip=8mm,
  sharp corners
}
\AtEndEnvironment{dem}{\null\hfill\textcolor{colorAsignatura!25!black}{$\blacksquare$}}

%====AFIRMACIONES====%
\newenvironment{claim}[1][Afirmación]{\par\noindent\textbf{#1: }\space\itshape}{\par}
\newenvironment{claimproof}[1][Demostración de la afirmación]{\par\noindent\textit{#1.}\space}{\hfill $\triangle$\par\medskip}

% ==========================================
% ENTORNO DE EJERCICIOS
% ==========================================
\NewTcbTheorem[number within=chapter]{ejer}{Ejercicio}
{label type=ejer,
    colback=colorAsignatura!3!white, enhanced, title=Ejercicio,
    attach boxed title to top left={xshift=-4mm}, boxrule=0pt,
    after skip=2mm, % <--- Salto pequeño para que la resolución se pegue
    before skip=1cm, right skip=0cm, breakable, fonttitle=\bfseries, 
    toprule=0pt, bottomrule=0pt, rightrule=0pt, leftrule=4pt, arc=0mm, skin=enhancedlast jigsaw, sharp corners, 
    colframe=colorAsignatura!25!black, colbacktitle=colorAsignatura!25!black,
    boxed title style={
        frame code={ 
            \fill[colorAsignatura!25!black]
                (frame.south west)--(frame.north west)--(frame.north east)--
                ([xshift=3mm]frame.east)--(frame.south east)--cycle;
            \draw[line width=1mm, colorAsignatura!25!black]
                ([xshift=2mm]frame.north east)--([xshift=5mm]frame.east)--([xshift=2mm]frame.south east);
            % Pliegue inferior izquierdo (sombra un poco más oscura)
            \fill[colorAsignatura!15!black]
                (frame.south west)--+(4mm,-2mm)--+(4mm,2mm)--cycle;
        }
    }
}{ejer}

% ==========================================
% ENTORNO DE RESOLUCIÓN (Concatenado)
% ==========================================
\newtcolorbox{resol}[1][Resolución]{
  enhanced,
  breakable,
  frame hidden, 
  boxrule=0pt,
  borderline west={4pt}{0pt}{colorAsignatura!25!black}, 
  colback=colorAsignatura!3!white, 
  coltitle=colorAsignatura!25!black, 
  fonttitle=\rmfamily\bfseries\scshape\large, 
  title={#1}, 
  before skip=0mm, 
  after skip=1cm,
  sharp corners
}
\AtEndEnvironment{resol}{\null\hfill\textcolor{colorAsignatura!25!black}{$\blacktriangleleft$}}


% ======================================================================= %
%                      11. USO DE CITAS Y REFERENCIAS      
% ======================================================================= %


% A. Aplicación de Epígrafe (justo debajo del capítulo)
% \epigraph{\textit{"El tiempo es la sustancia de la que estoy hecho."}}{--- Jorge Luis Borges}


% B. Citación inteligente con cleveref (detecta automáticamente que es un Axioma)
% \cref{ax:reflexividad} establece una restricción algebraica fundamental sobre el flujo temporal. A partir de este principio primitivo, evaluamos la verdad de las fórmulas mediante la relación de forzamiento o satisfactibilidad semántica, denotada por $\Vdash$.


% C. Declaración en el Glosario de Símbolos
%\nomenclature{$\Vdash$}{Relación de satisfactibilidad local o forzamiento en un modelo}

% D. Lista compacta gracias a enumitem (márgenes estrechos y numeración romana en cursiva)


% E. Añadido al Índice Alfabético Analítico
% \index{Relación de accesibilidad} 

% F. Cita Bibliográfica automatizada
% \cite{goldblatt1992logics}. 


% ==========================================
% PLANTILLA: TABLA DE VERDAD GENERALIZADA
% ==========================================
%\rowcolors{2}{white}{colorAsignatura!5} 
%\begin{table}[H]
%    \centering
%    \renewcommand{\arraystretch}{1.5} 
%    \setlength{\tabcolsep}{14pt} 
%    \begin{tabular}{ccc !{\color{colorAsignatura}\vrule width 1.5pt} ccc} 
%        \toprule
%        \rowcolor{colorAsignatura!15} 
%        % 1. Cambia las variables libres y los pasos intermedios aquí:
%        $p$ & $q$ & $r$ & \textit{Paso 1} & \textit{Paso 2} & \textbf{Fórmula Final} \\
%        \midrule
%        % 2. Rellena los valores a la derecha de la barra en cada fila:
%        \verd & \verd & \verd &  &  &  \\
%        \verd & \verd & \fals &  &  &  \\
%        \verd & \fals & \verd &  &  &  \\
%        \verd & \fals & \fals &  &  &  \\
%        \fals & \verd & \verd &  &  &  \\
%        \fals & \verd & \fals &  &  &  \\
%        \fals & \fals & \verd &  &  &  \\
%        \fals & \fals & \fals &  &  &  \\
%        \bottomrule
%    \end{tabular}
%    \caption{Evaluación semántica de la fórmula...}
%    \label{tab:verdad_generica}
%\end{table}
%\rowcolors{1}{}{} % Apaga el color alterno para el resto del texto

% ==========================================
% PLANTILLAS: REGLAS DE INFERENCIA
% ==========================================

%% Opción A: Inferencia con 1 premisa
%\begin{prooftree}
%    \AxiomC{$\alpha \land \beta$}
%    \RightLabel{\footnotesize (Eliminación $\land$)}
%    \UnaryInfC{$\alpha$}
%\end{prooftree}
%
%% Opción B: Inferencia con 2 premisas
%\begin{prooftree}
%    \AxiomC{$\alpha \to \beta$}
%    \AxiomC{$\alpha$}
%    \RightLabel{\footnotesize (Modus Ponens)}
%    \BinaryInfC{$\beta$}
%\end{prooftree}
%
%% Opción C: Inferencia con 3 premisas
%\begin{prooftree}
%    \AxiomC{$\alpha \lor \beta$}
%    \AxiomC{$\alpha \to \gamma$}
%    \AxiomC{$\beta \to \gamma$}
%    \RightLabel{\footnotesize (Eliminación $\lor$)}
%    \TrinaryInfC{$\gamma$}
%\end{prooftree}

% ==========================================
% PLANTILLA: DEDUCCIÓN SECUENCIAL FORMAL
% ==========================================
%\begin{dem}[Demostración de la Transitividad de la Implicación]
%Para demostrar formalmente el teorema $\vdash (p \to q) \land (q \to r) \to (p \to r)$, procedemos mediante la siguiente deducción secuencial:
%
%\vspace{0.3cm}
%\begin{center}
%    \renewcommand{\arraystretch}{1.4}
%    \begin{tabular}{rll}
%        % Estructura: Número. & Fórmula & Justificación \\
%        1. & $p \to q$ & (Premisa) \\
%        2. & $q \to r$ & (Premisa) \\
%        3. & $p$ & (Hipótesis auxiliar / Supuesto) \\
%        4. & $q$ & (Modus Ponens en 1 y 3) \\
%        5. & $r$ & (Modus Ponens en 2 y 4) \\
%        6. & $p \to r$ & (Teorema de la Deducción en 3 y 5) \\
%    \end{tabular}
%\end{center}
%\vspace{0.1cm}
%Por lo tanto, la fórmula queda deducida en el sistema axiomático.
%\end{dem}
\begin{document}
% ======================================================================= %
%                                  PORTADA                                %
% ======================================================================= %

\includepdf[pages=1]{Portada/portada.pdf}

% ======================================================================= %
%                           PÁGINA DE CORTESÍA                            %
% ======================================================================= %
\cleardoublepage 
\thispagestyle{empty} 
\vspace*{5cm}
\begin{center}
    \itshape
    "La lógica es la anatomía del pensamiento." \\
    \vspace{0.3cm}
    --- John Locke
\end{center}


% ======================================================================= %
%                 ÍNDICE GENERAL Y GLOSARIO DE SÍMBOLOS                   %
% ======================================================================= %
\frontmatter % Activa numeración romana y desactiva numeración de capítulos

\cleardoublepage 
\assignpagestyle{chapter}{empty} 
\pagestyle{empty} 

\tableofcontents

\cleardoublepage
\assignpagestyle{chapter}{plain} % Restauramos el estilo del capítulo
\pagestyle{fancy}

% ======================================================================= %
%                               PRELIMINARES                              %
% ======================================================================= %

\chapter{Conceptos Preliminares}

En este capítulo se recogen las definiciones topológicas y algebraicas que, si bien no pertenecen estrictamente al temario de la asignatura, resultan fundamentales para el desarrollo de las demostraciones posteriores, espeicialmente, los conceptos de filtro y ultrafiltro.

\section*{Preliminares axiomáticos}

En primer lugar, debemos establecer que trabajaremos con la axiomática clásica de  $\text{(Zermelo-Fraenkel)}$, complementada con el axioma de elección $\text{(AC)}$ para la teoría de conjuntos.

Por su reiterado uso a lo largo de la asignatura, vamos a exponer tres de las principales equivalencias del axioma de elección que usaremos:

\begin{defn*}{Equivalencias del Axioma de Elección}
El Axioma de Elección es equivalente a las siguientes afirmaciones:
\begin{enumerate}
    \item[(1)] Para cualquier conjunto $X$, existe una función de elección $\varphi: \mathcal{P}(X) \setminus \{\emptyset\} \to X$ tal que $\varphi(A) \in A$ para todo conjunto no vacío $A \subseteq X$.
    \item[(1')] Para cualquier conjunto $X$, para todo subconjunto $B \in \mathcal{P}(X)$, existe un elemento $x_{B} \in B$ y se cumple que $\{x_{B} : B \in \mathcal{P}(X)\} \subseteqq \mathcal{P}(X)$.
    \item[(2)] Todo conjunto no vacío admite un buen orden; esto es, que para cada conjunto $X$ existe una \textbf{\textit{relación de orden total}} $\prec$ sobre $X$ tal que todo subconjunto no vacío $N \subseteq X$ tiene un elemento mínimo respecto de $\prec$.
    \item[(3)] \textit{(Lema de Zorn)} Todo conjunto parcialmente ordenado no vacío $X$ en el que toda cadena (subconjunto totalmente ordenado) tiene un elemento superior, posee al menos un elemento maximal.
\end{enumerate}
\end{defn*}




\section*{Preliminares algebraicos}

Para esta sección, recordaremos nociones elementales del álgebra lineal

\begin{defn*}{Sistemas libres, generadores y bases de Hamel}
Dado un \textit{\textbf{espacio vectorial}} $V$ sobre un cuerpo $\mathbb{K} \in \{\RR,\CC\}$, decimos que un conjunto de vectores $\{v_{\ii}\}_{\ii \in I} \subseteq V$ es:
\begin{itemize}
    \item \textbf{Sistema Libre} si para cada $v \in V$ existe un subconjunto finito $J \subseteq I$ y para toda colección de esscalares $\{\lamb_{\jj}\}_{\jj \in J} \subseteq \KK$ se cumple que
    \[
        \sum_{\jj \in J} \lamb_{\jj} v_{\jj} = 0 \implies \lamb_{\jj} = 0, \quad \forall \jj \in J.
    \]
    \item \textbf{Sistema Generador} si para cada $v \in V$ existe un subconjunto finito $J \subseteq I$ y escalares $\{\lamb_{\jj}\}_{\jj \in J} \subseteq \KK$ tales que
    \[
        v = \sum_{\jj \in J} \lamb_{\jj} v_{\jj}.
    \]
    \item \textbf{Base de Hamel} si es a la vez libre y generador.
\end{itemize}

\end{defn*}

A partir de esto, probaremos un resultado clásico sobre la existencia de bases de Hamel para espacios vectoriales.

\begin{prop*}{Existencia de bases de Hamel}
Todo espacio vectorial sobre un cuerpo $\mathbb{K} $ admite una base de Hamel.
\end{prop*}

\begin{dem}
La prueba se basará en aplicar el \textit{Lema de Zorn} sobre un conjunto parcialmente ordenado de sistemas libres de un espacio vectorial.
\medskip
Sea V un espacio vectorial sobre un cuerpo $\mathbb{K}$. Consideremos el conjunto $\mathcal{C}$ de todos los sistemas libres de $V$. 
$$
\mathcal{C} \coloneq \{(v_{\ii})_{\ii \in I} \subseteq V  : \text{el conjunto } (v_{\ii})_{\ii \in I} \text{ es libre}\}
$$

Seguidamente, dotaremos a $\mathcal{C}$ de un orden parcial mediante la inclusión de conjuntos, de modo que $(\mathcal{C}, \subseteq)$ sea un conjunto parcialmente ordenado.
Consideamos una \textbf{cadena} $C \subseteq \mathcal{C}$ y veamos que admite una cota superior.
\medskip
Consideramos $\{v_{\ii}\}_{\ii \in I}$ el conjunto formado por la \textit{unión de todos los sitemas libres que forman parte de C}. Veamos que se trata de un sistema libre

Sea $\{\ii_{1}, \cdots , \ii_{\nn}\} \subseteq I$ un subconjunto finito con $\nn \in \NN$ y una colección de escalares $(\lamb_{\kk})_{\kk = 1}^{\nn} \subseteq \KK$. Por la definición de $\{v_{\ii}\}_{\ii \in I}$, existen una serie de \textit{sistemas libres} $\{B_{1}, \cdots, B_{\nn}\} \subseteq C$ que verifican
\[
\forall \kk \in \{1, \cdots, \nn\} \hspace{5pt} v_{\ii_{\kk}} \in B_{\kk}
\]


Ahora bien, por ser C cadena, debe existir un sistema libre $B_{\lal} \in C \lal \in \{1, \cdots, \nn\}$ que contiene a todos los demás. De este modo, se verifica que 

\[
\{v_{\ii_{1}}, \cdots, v_{\ii_{\nn}}\} \subseteq B_{\lal}
\]

De esto se deduce que si $\sum_{\kk = 1}^{\nn} \lamb_{\kk} v_{\ii_{\kk}} = 0$, entonces $\lamb_{\kk} = 0$ para todo $\kk \in \{1, \cdots, \nn\}$, por ser $B_{\lal}$ un sistema libre. Por lo tanto, $\{v_{\ii}\}_{\ii \in I} \in \mathcal C$.
Además, por construcción, $\{v_{\ii}\}_{\ii \in I}$ es una cota superior de C, ya que contine a todos sus elementos. Por el \textit{Lema de Zorn}, $\mathcal{C}$ posee un elemento maximal $\{w_{\jj}\}_{\jj \in J}$. Veamos que $\{w_{\jj}\}_{\jj \in J}$ es una base de Hamel de $V$. Por lo previamente demostrado, basta comprobar que se trata de un sistema generador.

\medskip
Supongamos por reducción al absurdo que no lo es, entonces existe un elemento $v \in V$ que no puede expresarse como combinación lineal finita de los elementos de $\{w_{\jj}\}_{\jj \in J}$, es decir, $v \in V \setminus span(\{w_{\jj}\}_{\jj \in J})$. Sin embargo, tenemos que $\{v\} \cup \{w_{\jj}\}_{\jj \in J}$ es un sistema libre.
\smallskip
En efecto, si tomamos una combinación lineal finita de los elementos;
\[
\lamb v + \sum_{\ii = 1}^{\nn} \lamb_{\jj} w_{\jj_{\ii}} = 0 \colon \{\jj_{\ii}\}_{\ii = 1}^{\nn} \subseteq J
\]

debemos distinguir dos casos:
\begin{itemize}
    \item Caso 1: $\lamb \neq 0$. En este caso, podemos despejar $v$ en términos de los $w_{\jj}$, es decir $v \in span(\{w_{\jj}\}_{\jj \in J})$, \textcolor{red}{\# lo que contradice la suposición de que $v$ no está en el span de $\{w_{j}\}_{j \in J}$.} 
    \item Caso 2: $\lamb = 0$. Entonces la combinación lineal se reduce a una combinación lineal de los elementos $\{w_{\jj}\}_{\jj \in J}$, y por ser este un sistema libre, todos los coeficientes deben ser cero, incluyendo los $\lamb_{\jj}$.
\end{itemize}

\smallskip
En cualquier caso, colcuímos que $\{v\} \cup \{w_{\jj}\}_{\jj \in J}$ es un sistema libre,  \textcolor{red}{\# lo que contradice la maximalidad de $\{w_{j}\}_{j \in J}$ en $\mathcal{C}$.} Por lo tanto, conlcuímos que $\{w_{\jj}\}_{\jj \in J}$ es una base de Hamel de $V$.
\end{dem}


De hecho, podemos afirmar un resultado algo más fuerte.

\begin{cor*}{Extensión de sistemas libres a bases de Hamel}
Dado un espacio vectorial $V$ sobre un cuerpo $\KK$ y un sistema libre suyo $(z_{\ii})_{\ii \in I}$, se puede conseguir una base de Hamel de  $V$ que contiene a todos los elementos de $(z_{\ii})_{\ii \in I}$.
\end{cor*}
\begin{dem}
Basta tomar $\mathcal C$ como el conjunto de todos los sistemas libres de $V$ que contienen a $(z_{\ii})_{\ii \in I}$ y aplicar el mismo razonamiento que en la demostración anterior.  
\end{dem}   


\section*{Preliminares de Topología}
Como matiz inicial, debemos aclarar que todos los espacios topológicos descritos, salvo que se especifique lo contrario, se considerarán \textit{Hausdorff} ($T_{2}$).

\medskip
En esta sección, abordaremos cuestiones fundamentales de compacidad en términos de \textit{conjuntos cerrados}

\begin{defn*}{Propiedad de Intesección Finita (PIF)}
    Una colección de subconjuntos $\{K_{\ii} : \ii \in I\}$ de un espacio topológico $(X, \taut)$ se dice que tiene la \textbf{Propiedad de Intersección Finita (PIF)} si para todo subconjunto finito $J \subseteq I$ se cumple que
    \[
    \bigcap_{\jj \in J} K_{\jj} \neq \emptyset
    \]
\end{defn*}

Esta propiedad es fundamental para caracterizar la compacidad en términos de conjuntos cerrados, como veremos a continuación.

\begin{prop*}{Caracterización de la compacidad mediante la PIF}
   Dado un espacio topológico $(X, \taut)$, los siguientes enunciados son equivalentes:
    \begin{enumerate}
          \item[(1)] $(X, \taut)$ es compacto.
          \item[(2)] Para toda colección de conjuntos cerrados $\{K_{\ii} : \ii \in I\}$ que posee la PIF, se cumple que
          \[
          \bigcap_{\ii \in I} K_{\ii} \neq \emptyset
          \]
     \end{enumerate}
\end{prop*}

\begin{dem}
    Realizaremos la demostración de cada una de las direcciones

    \medskip
    $\borde{(1) \implies (2)}$ Supongamos que $(X, \taut)$ es compacto y consideremos una colección de conjuntos cerrados $\{K_{\ii} : \ii \in I\}$ cumpliendo que $\bigcap_{\ii \in I} K_{\ii} = \emptyset$. Entonces , por las propiedades de la complementación se tiene que
    \[
    \begin{aligned}
        X = X \setminus \emptyset &= X \setminus \bigcap_{\ii \in I} K_{\ii} = \bigcup_{\ii \in I} (X \setminus K_{\ii})
    \end{aligned}
    \]
    donde $\{X \setminus K_{\ii} : \ii \in I\}$ es una colección de conjuntos abiertos que cubre $X$. Por haber supuesto que $(X, \taut)$ es compacto, existe un subconjunto finito $J \subseteq I$ tal que $\{X \setminus K_{\jj} : \jj \in J\}$ cubre $X$.
    Entonces, nuevamente por las propiedades de la complementación, se tiene que:

    \[
    \begin{aligned}
        X &= \bigcup_{\jj \in J} (X \setminus K_{\jj}) = X \setminus \bigcap_{\jj \in J} K_{\jj}\\
         &\implies \emptyset= \bigcap_{\jj \in J} K_{\jj}
    \end{aligned}
    \]

    Luego, se concluye que $\{K_{\ii} : \ii \in I\}$ no posee la PIF.

    \medskip
    $\borde{(2) \implies (1)}$ Supongamos que toda colección de conjuntos cerrados $\{K_{\ii} : \ii \in I\}$ que posee la PIF cumple que $\bigcap_{\ii \in I} K_{\ii} \neq \emptyset$. Consideremos un recubrimiento abierto arbitrario de $X$,  $\{V_{\ii} : \ii \in I\}$. Entonces, se tiene que:

    \[
    \begin{aligned}
        X &= \bigcup_{\ii \in I} V_{\ii} \\
        &\implies \emptyset = X \setminus \bigcup_{\ii \in I} V_{\ii} = \bigcap_{\ii \in I} (X \setminus V_{\ii}) 
    \end{aligned}
    \]
    \medskip

    Como $\{X \setminus V_{\ii} : \ii \in I\}$ es una colección de conjuntos cerrados cumpliendo que $\bigcap_{\ii \in I} (X \setminus V_{\ii}) = \emptyset$, se concluye por hipótesis  que $\{X \setminus V_{\ii} : \ii \in I\}$ no posee la PIF. Por lo tanto, existe un subconjunto finito $J \subseteq I$ tal que
    \[
    \begin{aligned}
     \emptyset &= \bigcap_{\jj \in J} (X \setminus V_{\jj}) = X \setminus \bigcup_{\jj \in J} V_{\jj}\\
     & \implies X = \bigcup_{\jj \in J} V_{\jj}
    \end{aligned}
    \]

    Luego, por definición, $(X, \taut)$ es compacto.

\end{dem}



\section*{Filtros y Ultrafiltros}

\subsection*{Filtros y ultrafiltros sobre conjuntos}

\begin{defn*}{Filtros sobre conjuntos}
Dado un conjunto infinito $I$ (evidentemente no vacío), a cuyos elementos denominaremos \textit{índices} y una colección de subconjuntos $\filt$ de $I$, decimos que $\filt$ es un \textbf{filtro} sobre $I$ si cumple las siguientes propiedades:
\begin{enumerate}
    \item[(i)] $\emptyset \notin \filt$ y $I \in \filt$.
    \item[(ii)] Si $A, B \in \filt$, entonces $A \cap B \in \filt$.
    \item[(iii)] Si $A \in \filt$ y $A \subseteq B \subseteq I$, entonces $B \in \filt$.
\end{enumerate}
\end{defn*}

Consideramos ahora algunos ejemplos fundamentales de filtros sobre conjuntos.

\begin{ejemplo*}{Filtro trivial}
El conjunto $\filt = \{I\}$ es un filtro sobre $I$, denominado \textit{filtro trivial}.

\medskip
En efecto, se verifica que $\emptyset \notin \filt$ y $I \in \filt$. Además, si $A, B \in \filt$, entonces $A = B = I$, por lo que $A \cap B = I \in \filt$. Finalmente, si $A \in \filt$ y $A \subseteq B \subseteq I$, entonces $B = I$, por lo que $B \in \filt$.

Este ejemplo, permite asegurar la existencia de filtros sobre cualquier conjunto infinito.
\end{ejemplo*}

\begin{ejemplo*}{Filtros asociados a un elemento}
Sea $i_{0} \in I$ un elemento fijo. Consideremos el conjunto $\filt_{i_{0}} = \{A \subseteq I : i_{0} \in A\}$. Entonces, $\filt_{i_{0}}$ es un filtro sobre $I$, denominado \textit{filtro principal asociado al elemento $i_{0}$}.
    
\medskip
En efecto, dado un índice $i_{0} \in I$, como estamos tomando subconjuntos de $I$ que lo contienen, se cumple por las propiedades de conjuntos que $\emptyset \notin \filt_{i_{0}}$ y $I \in \filt_{i_{0}}$. Además, si $A, B \in \filt_{i_{0}}$, entonces $i_{0} \in A$ y $i_{0} \in B$, por lo que $i_{0} \in A \cap B$, es decir, $A \cap B \in \filt_{i_{0}}$. Finalmente, si $A \in \filt_{i_{0}}$ y $A \subseteq B \subseteq I$, entonces $i_{0} \in A$ y por lo tanto $i_{0} \in B$, es decir, $B \in \filt_{i_{0}}$.

\end{ejemplo*}



\begin{ejemplo*}{Filtro de Fréchet}
    Sea $I$ un conjunto infinito. Consideremos el conjunto $\filt = \{A \subseteq I : I \setminus A \text{ es finito}\}$. Entonces, $\filt$ es un filtro sobre $I$, denominado \textit{filtro de Fréchet}.

    \medskip
    En efecto, vamos a verificar las propiedades de un filtro:
    \begin{enumerate}
        \item[(i)] Como $I$ es infinito y $\emptyset = I \setminus I$, se sigue que $\emptyset \not \in \filt$. Por otro lado, como $\emptyset$ es finito y $I = I \setminus \emptyset$, se sique que $I \in \filt$.
        \item[(ii)] Sean $A, B \in \filt$. Entonces, $I \setminus A$ y $I \setminus B$ son conjuntos finitos. Por lo tanto, su unión $(I \setminus A) \cup (I \setminus B)$ también es finita. Además, por las propiedades de De Morgan, se tiene que
        \[
        I \setminus (A \cap B) = (I \setminus A) \cup (I \setminus B)
        \]
        lo que implica que $A \cap B \in \filt$.
        \item[(iii)] Sea $A \in \filt$ y $A \subseteq B \subseteq I$. Entonces, $I \setminus A$ es finito. Además, por las propiedades de conjuntos, se tiene que
        \[
        (I \setminus B) \subseteq (I \setminus A)
        \]
        lo que implica que $I \setminus B$ es finito, y por tanto $B \in \filt$.
    \end{enumerate}
\end{ejemplo*}

\begin{prop*}{Bases y subbases de ultrafiltros}
Sea $I$ un conjunto infinito y sea $\mathcal S \subset \mathcal P(I)$ una colección de subconjuntos con la PIF. Entonces, la familia de sunconjuntos:
$$
\filt_{\mathcal{S}} \coloneq \left\{ A \subseteq I \colon A \supseteq \bigcap_{F \in \mathcal{S}'} F \quad\text{con}\quad \mathcal{S}' \subseteq \mathcal{S} \text{ finito} \right\}
$$

es un filtro sobre $I$ y además, es el filtro sobre $I$ más pequeño que continene a $\mathcal S$. Es decir, si \textcolor{magenta}{$\mathcal G$} es un filtro sobre $I$ que contiene a $\mathcal S$, entonces $\filt_{\mathcal{S}} \subseteq \textcolor{magenta}{\mathcal G}$.

\end{prop*}


\begin{dem}

    Vamos a verificar que $\filt_{\mathcal{S}}$ cumple las propiedades de un filtro:
    \begin{enumerate}
        \item[(i)] Notemos que si $\emptyset \in \filt$, entonces existe una subcolección finita $\mathcal{S}' \subseteq \mathcal{S}$ tal que $\emptyset \supseteq \bigcap_{F \in \mathcal S'} F$ \textcolor{red}{\# esto contradice el hecho de que $\mathcal S$ tiene la PIF}. Por lo tanto, $\emptyset \notin \filt_{\mathcal{S}}$.Deduciremos que $I \in \filt_{\mathcal{S}}$ a partir de \emph{(iii)}
        \item[(ii)] Sean $A, B \in \filt_{\mathcal S}$. Entonces, existen subcolecciones finitas $\mathcal{S}_{A}, \mathcal{S}_{B} \in \mathcal S$tal que:
    \[
     \left.
     \begin{aligned}
            A &\supseteq \bigcap_{F \in \ \mathcal{S}_{A}} F \\
            B &\supseteq \bigcap_{F \in \ \mathcal{S}_{B}} F
        \end{aligned}
        \right\} \implies A \cap B \supseteq \bigcap_{F \in \mathcal{S}_{A} \cup \mathcal{S}_{B}} F \colon \mathcal{S}_{A} \cup \mathcal{S}_{B} \subseteq \mathcal S' \text{ es finita}
       \]
       De esto se deduce que $A \cap B \in \filt_{\mathcal S}$.
        \item[(iii)] Sea $A \in \filt_{\mathcal S}$ y $A \subseteq B \subseteq I$. Entonces, existe una subcolección finita $\mathcal{S}_{A} \subseteq \mathcal S$ tal que $A \supseteq \bigcap_{F \in \mathcal{S}'} F$. Por lo tanto, se tiene que:
        \[
        B \supseteq A \supseteq \bigcap_{F \in \mathcal{S}'} F \implies B \in \filt_{\mathcal S}
        \]
    \end{enumerate}

    \medskip
    Sea \textcolor{magenta}{$\mathcal G$} un filtro sobre $I$ que contiene a $\mathcal S$. Entonces, por la propiedad (ii) de filtro, para cada subfamilia finita $\mathcal S' \subseteq \mathcal S$ se cumple que
    \(\bigcap_{F \in \mathcal S'} F \in \textcolor{magenta}{\mathcal G}\) y por la propiedad (iii) de filtro, para cada $A \supset \bigcap_{F \in \mathcal S'} F$ se cumple que $A \in \textcolor{magenta}{\mathcal G}$. Por lo tanto, para cada $A \in \filt_{\mathcal S}$ se tiene que $A \in \textcolor{magenta}{\mathcal G}$. Es decir, $\filt_{\mathcal S} \subseteq \textcolor{magenta}{\mathcal G}$.
    
\end{dem}

Este resultado nos lleva a definir algunas nociones que resultan particularmente intuitivas al tratar con estructuras de esta naturaleza:

\begin{defn*}{Filtros generados, bases y subbases de filtros}
    Bajo las condiciones del resultado anterior, definimos:
    \begin{itemize}
        \item \textbf{Filtro generado por una familia de conjuntos:} Sea $\mathcal S$ una familia de conjuntos no vacíos. El filtro generado por $\mathcal S$, denotado por $\filt_{\mathcal S}$, es el menor filtro que contiene a $\mathcal S$.
        \item \textbf{Subbase de un filtro:} Una subbase de un filtro $\mathcal F_{\mathcal S}$ es una familia no vacía de subconjuntos que cumple la PIF y lo genera.
        \item \textbf{Base de un filtro:} Una base de un filtro $\mathcal F_{\mathcal S}$  es una familia no vacía de subconjuntos que verifica que toda intersección finita de sus elementos pertenece a sí misma y lo genera. Es decir, si para cada subfamilia finita 
        $\mathcal S' \subseteq \mathcal S$ entonces $\bigcap_{F \in \mathcal S'} F \in \mathcal S$.
    \end{itemize}
\end{defn*}

Ilustremos estas definiciones con un ejemplo sencillo sobre los números naturales.

\begin{ejemplo*}{El filtro de Fréchet sobre \NN}
    Denotemos por $\filt$ al filtro de Fréchet sobre $\NN$, es decir:
    \[
    \filt = \{A \subseteq \NN : \NN \setminus A \text{ es finito}\}
    \]

    Veamos un ejemplo de subbase y base para este filtro:
    \begin{itemize}
        \item \textbf{Subbase:} Consideremos la familia de conjuntos $\mathcal S = \{\NN \setminus \{\nn\} : \nn \in \NN\}$. Esta familia cumple la PIF.

        \medskip
        En efecto, para cualquier subconjunto finito $J \subseteq \NN$, la intersección de los conjuntos $\NN \setminus \{\nn\}$ para $\nn \in J$ es:
        \[
        \bigcap_{n \in J} \  (\NN \setminus \{\nn\}) = \NN \setminus J \neq \emptyset
        \]
        que es un conjunto cofinito.
        
        \medskip
        Veamos que en efecto, $\mathcal S$ genera el filtro de Fréchet, es decir, $\filt_{\mathcal S} = \filt$.

        \medskip
        \borde{\subseteq} Tomamos un elemento del filtro generado por $\mathcal S$, es decir, un conjunto $A \in \filt_{\mathcal S}$. Por definición, existe una subfamilia finita $\mathcal S' \subseteq \mathcal S$ tal que:
        \[
        A \supseteq \bigcap_{F \in \mathcal S'} F
        \]
        Como cada $F \in \mathcal S'$ es un conjunto cofinito, la intersección finita de estos conjuntos también es un conjunto cofinito. Por lo tanto, $A$ es un conjunto cofinito, lo que implica que $A \in \filt$. Así, $\filt_{\mathcal S} \subseteq \filt$.

        \smallskip
        \borde{\supseteq} Ahora, tomamos un elemento del filtro de Fréchet, es decir, un conjunto $A \in \filt$. Por definición, $\NN \setminus A$ es finito. Consideremos la familia $\mathcal S' = \{\NN \setminus \{\nn\} : \nn \in \NN \setminus A\}$. Esta familia está contenida en $\mathcal S$ y su intersección es:
        \[
        \bigcap_{\nn \in \NN \setminus A} (\NN \setminus \{\nn\}) = A
        \]
        Por lo tanto, $A$ pertenece al filtro generado por $\mathcal S$, es decir, $A \in \filt_{\mathcal S}$. Así, $\filt \subseteq \filt_{\mathcal S}$.

        En resumen, $\filt_{\mathcal S} = \filt$.

        \medskip
        \item \textbf{Base:} Consideremos la familia de conjuntos $\mathcal B = \{\NN \setminus \{1, 2, \ldots, \nn\} : \nn \in \NN\}$. Veamos en primer lugar, que el conjunto cumple la PIF.
        
        \medskip
        Si tomamos una colección finita de subconjuntos $\{\NN \setminus \{1, \cdots, \ii_{\kk}\}\}_{\kk=1}^{\nn}$, \ por ser $(\NN,\leq)$ un \textbf{orden total}
        , podemos asumir sin perdida de generalidad que los elementos satisfacen $\forall \kk \in \{1, \cdots, \nn -1 \} \ \ii_{\kk} < \ii_{\kk+1}$. De este modo, se tiene que:
        \[
        \bigcap_{\kk = 1}^{\nn} \ (\NN \setminus \{1, \cdots \ii_{\kk}\}) = \NN \setminus \{1, \cdots, \ii_{\nn}\} \in \mathcal B
        \]

        Por lo tanto, como cumple la condición adicional de base, basta ver que genera el filtro de Fréchet, nuevamente, que $\filt_{\mathcal B} = \filt$.

        \medskip
        \borde{\subseteq} Tomamos un elemento $A \in \filt_{\mathcal B}$. Por definición, existe una subfamilia finita $\mathcal B' \subseteq \mathcal B$ tal que:
        \[
            A \supseteq \bigcap_{F \in \mathcal B'} F
        \]
        Como cada $F \in \mathcal B'$ es un conjunto cofinito, la intersección finita de estos conjuntos también es un conjunto cofinito (por lo visto anteriormente). Por lo tanto, $A$ es un conjunto cofinito, lo que implica que $A \in \filt$. Así, $\filt_{\mathcal B} \subseteq \filt$.

        \smallskip
        \borde{\supseteq} Ahora, tomamos un elemento $A \in \filt$. Por definición, $\NN \setminus A$ es finito. Sea $\nn = \max(\NN \setminus A)$. Entonces, si consideramos el conjunto $\NN \setminus \{1, 2, \ldots, \nn\} \in \mathcal B$, se tiene que:
        \[
            A \supseteq \NN \setminus \{1, 2, \ldots, \nn\} = \bigcap_{\ii = 1}^{\nn} \ (\NN \setminus \{1, 2, \ldots, \ii\}) 
        \]
        Por lo tanto, $A \in \filt_{\mathcal B}$. Así, $\filt \subseteq \filt_{\mathcal B}$.
        

        En resumen, $\filt_{\mathcal B} = \filt$.

        \medskip
    \end{itemize}
\end{ejemplo*}
    
\medskip
Veamos ahora un resultado estándar para este tipo de estrcutras:la existencia de elementos maximales.

\begin{prop*}{Exstencia de ultrafiltros}
    Sean $I$ un conjunto infinito y $\filt$ un filtro sobre $I$. Entonces, existe un \textit{filtro maximal sobre $I$}, $\ultf$, que contiene a $\filt$.
\end{prop*}

\medskip
\begin{dem}
    La demostración se basa en la aplicación del \textit{Lema de Zorn} sobre un conjunto parcialmente ordenado de filtros sobre $I$.

    \medskip
    Consideramos el conjunto de filtros que contienen a $\filt$ como:
    \[
        \mathbb A = \{\gilt \subseteq \mathcal P(I) \mid \gilt \text{ es un filtro y } \filt \subseteq \gilt\}   
    \]
    al que dotamos además de una estrcutra de orden parcial a través de la relación de inclusión de conjuntos, $(\mathbb A, \subseteq)$.

    \smallskip
    Dada una cadena $C \subseteq \mathbb A$, consideremos el conjunto $\mathcal D = \bigcup_{\gilt \in C} \gilt$, que verifica lo siguiente:

    \begin{itemize}
        \item[$\bullet$]$\mathcal D$ es un filtro.
        \begin{enumerate}
            \item[(i)] Si suponemos por reducción al absurdo que $\emptyset \in \mathcal D$, entonces debe existir un filtro $\gilt \in C$ tal que $\emptyset \in \gilt$ \textcolor{red}{\# lo que contradice la definición de filtro}. Por otro lado, como $\forall \gilt \in C, I \in \gilt$, se sigue que $I \in \mathcal D$
            \item[(ii)] Sean $A, B \in \mathcal D$. Entonces, existen filtros $\gilt_{A}, \gilt_{B} \in C$ tales que $A \in \gilt_{A}$ y $B \in \gilt_{B}$. Como $C$ es una cadena, o bien $\gilt_{A} \subseteq \gilt_{B}$ o $\gilt_{B} \subseteq \gilt_{A}$. Supongamos que $\gilt_{A} \subseteq \gilt_{B}$. Entonces, $A \in \gilt_{B}$ y $B \in \gilt_{B}$, por lo que $A \cap B \in \gilt_{B}$. Por lo tanto, $A \cap B \in \mathcal D$. El otro caso es completamente análogo.
            \item[(iii)] Sea $A \in \mathcal D$ y $A \subseteq B \subseteq I$. Entonces, existe un filtro $\gilt_{A} \in C$ tal que $A \in \gilt_{A}$. Como $\gilt_{A}$ es un filtro y $A \subseteq B$, se tiene que $B \in \gilt_{A}$. Por lo tanto, $B \in \mathcal D$.
        \end{enumerate}

        \smallskip
        \item[$\bullet$] Por construcción, es claro que $\filt \subseteq \mathcal D$, por lo que $\mathcal D \in \mathbb A$. Además, $\mathcal D$ es cota superior de $C$, ya que para cada $\gilt \in C$, $\gilt \in \mathcal D$ 
    \end{itemize}

    \medskip
    Por el \textit{Lema de Zorn}, $\mathbb A$ posee un elemento maximal $\ultf$, que por construcción contiene a $\filt$.
\end{dem}

\medskip
Este resultado justifica la siguiente definición:

\begin{defn*}{Ultrafiltros}
    Un \textbf{ultrafiltro} sobre un conjunto infinito $I$ es un filtro maximal sobre $I$.
\end{defn*}

\smallskip
A pesar de tener garantizada la existencia de ultrafiltros, su definición es algo confusa en la práctica. Por ello, introduciremos una caracterización que nos permitirá abordar el concepto con mayor facilidad.

\begin{prop*}{Caracterización de ultrafiltros}
    Sea $I$ un conjunto infinito y $\filt$ un filtro sobre $I$. Entonces, los siguientes enunciados son equivalentes
     \begin{enumerate}
        \item[(1)] $\filt$ es un \textit{ultrafiltro}
        \item[(2)] Para cada sunconjunto $A \subseteq I$, o bien $A \in \filt$ o bien $I \setminus A \in \filt$
     \end{enumerate}
\end{prop*}

     \begin{dem}
     Veamos cada una de las implicaciones por separado

    \medskip
    \borde{(1) \implies (2)} \medskip Sea $A \subseteq I$ cualquiera. Supongamos por reducción al absurdo que $A \notin \filt$ y $I \setminus A \notin \filt$.
    
    \smallskip
    Entonces, como $A \notin \filt$, debe existir un subconjunto $F_{A} \in \filt$ tal que $A \cap F_{A} = \emptyset$. En efecto, si para cada $F \in \filt$ $A \cap F \neq \emptyset$, entonces $\filt \cup \{A\}$ tiene la PIF. De esta forma, por la definición de filtro generado, podemos considerar el filtro generado por $\filt \cup \{A\}$, $\gilt_{\filt \cup \{A\}}$, que por definición, cumple que
    \[
    \filt \subseteq \filt \cup \{A\} \subseteq \gilt_{\filt \cup \{A\}}. \textcolor{red}{\#}
    \]

    \smallskip
    Razonando de forma análoga para $B \coloneq I \setminus A$, deducirmos que existe cierto subconjunto $F_{B} \in \filt$, verificando que $B \cap F_{B} = \emptyset$. De ambas deducciones, se tiene que

    \[
    \left.
    \begin{aligned}
        F_{A} \cap A = \emptyset& \implies F_{A} \subset I \setminus A = B\\
        F_{B} \cap B  = \emptyset& \implies F_{B} \subset I \setminus B = A
    \end{aligned}
    \right\} \emptyset = F_{A} \cap F_{B} \in \filt \textcolor{red}{\#}
    \]

    \medskip

    \borde{(2) \implies (1)} Supongamos que $\filt$ es un filtro no maximal; ergo, por el resultado previo, existe un ultrafiltro $\ultf$ tal que $\filt \in \ultf$.

    \smallskip
    De esta forma, debe de existir un sunconjunto $A \subseteq I$ tal que $A \in \ultf \setminus \filt$. Por ser $\ultf$ filtro $I \setminus A \not \in \ultf$ (ya que de lo contrario $\emptyset \in \ultf$). Por lo tanto, $I \setminus A \not \in \filt$ y $A \not \in \filt$, lo que niega $(2)$.
     \end{dem}

     \subsection*{Filtros, ultrafiltros y topología}

     Esta sección pretende introducir una generalización del concepto de sucesiones y convergencia en espacios méttricos para cualquier espacio topológico.

     \begin{defn*}{Convergencia de filtros}
     Sea $(X, \taut)$ un espacio topológico y $\filt$ un filtro o ultrafiltro sobre un conjunto de índice $I$. Decimos que una familia de elementos $(x_{\ii})_{\ii \in I} \subseteq X$ \textbf{converge} a un elemento $x_{0} \in X$ según el filtro $\filt$ si:
    \begin{center}
        \borde{\forall V \in \mathcal{N}_{x_{0}}  \ \{\ii \in I \colon x_{\ii} \in V\} \in \filt}
    \end{center}
     
     \end{defn*}

     \medskip
     Ilustremos esta definición via un ejemplo sencillo en un espacio métrico.

     \begin{ejemplo*}{Convergencia de sucesiones según el filtro de Fréchet}
        Si consideramos el espacio topológico $(\RR, \taut_{d})$ de la recta real con su topología usual y el filtro de Fréchet sobre $\NN$, entonces la noción tradicional de convergencia coincide con la convergencia según el filrtro de Fréchet.

        \smallskip
        Veamos ambas direcciones de la equivalencia  para cerciorarnos de la equivalencia.
        \begin{itemize}
            \item[$\bullet$] Tomemos una sucesión $(x_{\nn})_{\nn \in \NN} \subseteq \RR$ y un elemento $x_{0} \in \RR$ tal que $x_{\nn} \underset{\nn}{x_{0}}$. Por definición, tenemos que:
             $$\forall \eps > 0 \ \exists \nn_{0} \in \NN \colon \forall \nn \geq \nn_{0} \implies |x_{\nn} - x_{0}| < \eps$$

             Sin embargo, basta notar que esto puede reformularse como sigue:
             $$
             \{\nn \in \NN \colon x_{\nn} \in (x_{0} - \eps, x_{0} + \eps)\} \supset \{\nn \in \NN \colon \nn \geq \nn_{0}\} \in \filt
             $$
             Donde la última inclusión es clara por ser el conjunto cofinito. Por las propiedades de filtro, se sigue que $\{\nn \in \NN \colon x_{\nn} \in (x_{0} - \eps, x_{0} + \eps)\} \in \filt$. Por lo tanto, $(x_{\nn})_{\nn \in \NN}$ converge a $x_{0}$ según el filtro de Fréchet.

             \item[$\bullet$] Tomemos ahora una sucesión $(x_{\nn})_{\nn \in \NN} \subseteq \RR$ y un elemento $x_{0} \in \RR$ tal que $(x_{\nn})_{\nn \in \NN}$ converge a $x_{0}$ según el filtro de Fréchet. Por definición, tenemos que:
             $$
             \forall \eps > 0 \ I_{\eps} \coloneq \{\nn \in \NN \colon x_{\nn} \in (x_{0} - \eps, x_{0} + \eps)\} \in \filt
             $$
             Luego, $I_{\eps}$ es cofinito; es decir, si tomamos la base del filtro de Fréchet, existe una colección de naturales $\{\kk_{1}, \cdots , \kk_{\mm}\} \colon \mm \in \NN$ verificando que:
             $$
                I_{\eps}  = \NN \setminus \{\kk_{1}, \cdots, \kk_{\mm}\}
             $$
             De esta forma, tomando $\kk = max\{\kk_{1}, \cdots, \kk_{\mm}\}$, tenemos que 
             
             $$
             \NN \setminus \{\kk_{1}, \cdots, \kk_{\mm}\} \supseteq \NN \setminus \{1, \cdots, \kk\} 
             $$.

             O expresado de otro modo, para cualquier $\nn \geq \kk$ se cumple que $x_{\nn} \in I_{\eps}$. Por lo tanto, $(x_{\nn})_{\nn \in \NN}$ converge a $x_{0}$ en el sentido tradicional.
        \end{itemize}
     \end{ejemplo*}

     \medskip

     La notación para denotar la convergencia respecto a un filtro suele ser bastante variada, aunque preferiblemente emplearemos la siguiente notación:

     \[
     \begin{array}{llllc}
        x_{\ii} \underset{\ii \to \filt}{\overset{\taut}{\to}} x_{0} & x_{\ii} \underset{\filt}{\overset{\taut}{\to}} x_{0} & \underset{\ii \to \filt}{Lim} x_{\ii} =  x_{0} 
     \end{array}
     \]

     \medskip

     Ahora que hemos establecido esta noción general de convergencia, vamos a introducir los elementos necesarios para abordar el principal caso que vamos a manejar.

     \begin{defn*}{Conjunto dirigidos}
        Decimos que una tuplam$(I, \preceq)$ es un \textbf{conjunto dirigido} si $I$ es un conjunto infinito y $\preceq$ es una relación de orden parcial sobre $I$ que verifica la siguiente condición:
         
        \begin{center}
            $
            \forall \ii, \jj \in I \ \exists \kk \in I : \ii \preceq \kk \text{ y } \jj \preceq \kk
            $
        \end{center}

        A esta relación la denominamos \textbf{relación de "mejor que"}
     \end{defn*}

     \smallskip
     A partir de un conjunto dirigido, $(I,\preceq)$, denotamos para cada índice $\ii \in I$
        \[
            I(\ii) \coloneq \{\jj \in I : \ii \preceq \jj\}    
        \]

    \smallskip
    Así, la colección $\{I(\ii) \colon \ii \in I\}$ tiene la PIF.

    \medskip
    En efecto, si tomamos $\ii, \jj \in I$, por ser $(I,\preceq)$ un conjunto dirigido, existe $\kk \in I$ tal que $\ii \preceq \kk$ y $\jj \preceq \kk$. Por lo tanto, $\kk \in I(\ii) \cap I(\jj)$, y por lo tanto $I(\ii) \cap I(\jj) \neq \emptyset$.

    \smallskip
    De esta forma, podemos considerar el filtro generado por la colección $\{I(\ii) \colon \ii \in I\}$, al que llamaremos \textbf{filtro de orden sobre $I$}

    \begin{rem*}{}
    Notemos que para cada $\ii \in I$, si $\jj \in I(\ii)$, entonces $I(\jj) \subseteq I(\ii)$. Basta notar que, si $\kk \in I(\jj)$, entonces $\jj \preceq \kk$. Como $\ii \preceq \jj$, por la transitividad de la relación de orden, se tiene que $\ii \preceq \kk$, es decir, $\kk \in I(\ii)$.
    \end{rem*}

    \begin{ejemplo*}{Caso fundamental}
        Sean $(X, \taut)$ un espacio topológico y $x_{0} \in X$. Denotando por $\mathcal{N}_{x_0}$ al conjunto de todos los entornos abiertos de $x_{0}$, si lo dotamos de la relación $A \preceq B \Leftrightarrow B \subseteq A$, entonces $(\mathcal{N}_{x_0}, \preceq)$ es un conjunto dirigido.
        
        \smallskip
        En efecto, si tomamos $A, B \in \mathcal{N}_{x_0}$, por la definición de entorno abierto, existe un entorno abierto $C$ de $x_{0}$ tal que $C \subseteq A \cap B$. Por lo tanto, $C \preceq A$ y $C \preceq B$, y por lo tanto $(\mathcal{N}_{x_0}, \preceq)$ es un conjunto dirigido.

        \smallskip
        Además, como $\mathcal{N}_{x_0}$ cumple la PIF, podemos considerar el filtro de orden sobre $\mathcal{N}_{x_0}$ respecto de $\preceq$, que denotaremos $\filt$

        \medskip
        Tomemos ahora un filtro sobre $I$, $\gilt$ que contiene a $\filt$. Entonces, para cada $V \in \mathcal{N}_{x_0}$, se cumple que
        \[
        x_{V} \underset{V \to \gilt}{\overset{\taut}{\to}} x_{0}
        \]

        \medskip
        Veamos que en efecto esto es cierto. Dado $W \in \mathcal{N}_{x_0}$, como si $V \preceq W$ entonces $x_{V} \in V \subseteq W$, deducimos que:
        \[
        \begin{aligned}
            \{V \in \mathcal{N}_{x_0} : x_{V} \in W\} \supseteq \{W\} &\implies \{V \in \mathcal{N}_{x_0} : x_{V} \in W\} \supseteq I(W) \\
            &\implies \{V \in \mathcal{N}_{x_0} : x_{V} \in W\} \in \gilt
        \end{aligned}
        \]
    \end{ejemplo*}

    Esto nos permite emular la caracterización sucesional de la clausura de un conjunto en términos de la convergencia respecto de un filtro.

    \begin{prop*}{Caracterización de la clausura en términos de filtros}
        Sea $(X, \taut)$ un espacio topológico y $A \subseteq X$. Entonces, los siguientes enunciados son equivalentes:

        \begin{enumerate}
            \item[(1)] $x_{0} \in \overline{A}$
            \item[(2)] Existen un conjunto infinito $I$, un filtro sobre $I$, $\filt$ y una colección de elementos $(x_{\ii})_{\ii \in I} \subseteq A$ tal que $x_{\ii} \underset{\ii \to \filt}{\overset{\taut}{\to}} x_{0}$.
            \item[(3)] Existen un conjunto infinito $I$, un ultrafiltro sobre $I$, $\ultf$ y una colección de elementos $(x_{\ii})_{\ii \in I} \subseteq A$ tal que $x_{\ii} \underset{\ii \to \ultf}{\overset{\taut}{\to}} x_{0}$.
        \end{enumerate}
    \end{prop*}

    \begin{dem}
        \borde{(1) \implies (2)} Sea $x_{0} \in \overline{A}$. Consideremos el conjunto dirigido $(\mathcal{N}_{x_0}, \preceq)$, el filtro de orden sobre $\mathcal{N}_{x_0}$, $\filt$ y para cada $\VV \in \mathcal{N}_{x_0}$, tomamos un elemento $x_{V} \in A \cap V$ (que existe por la hipótesis de que $x_{0} \in \overline{A}$). Entonces, por la observación del ejemplo previamente planteado:
        \[
        x_{V} \underset{V \to \filt}{\overset{\taut}{\to}} x_{0}
        \]

        \medskip
        \borde{(2) \implies (3)} Bajo las consideraciones de la implicación previa, basta aplicar la misma observación sobre el ultrafiltro $\ultf$.


        \medskip
        \borde{(3) \implies (1)} Sean un conjunto de índices $I$, un ultrafiltro sobre $I$ $\ultf$ y una colección de elementos $(x_{\ii})_{\ii \in I} \subseteq A$ tal que $x_{\ii} \underset{\ii \to \ultf}{\overset{\taut}{\to}} x_{0}$. Sea $V \in \mathcal{N}_{x_0}$ y veamos que $V \cap A \neq \emptyset$. Para ello, basta recordar la definición de convergencia en filtro, que asegura que:
        \[
        I_{V} = \{\ii \in I : x_{\ii} \in V\} \in \ultf
        \]
        Por ser $\ultf$ un ultrafiltro, sabemos que $I_{V} \neq \emptyset$, luego existe un índice $\jj \in I_{V}$. De este modo, por construcción, $x_{\jj} \in V \cap A$.
    \end{dem}


    \begin{prop*}{Caracterización de la continuidad en términos de filtros}
     Sean $(X,\taut_{X})$, $(Y,\taut_{Y})$ dos espacios topológicos, $f \colon (X, \taut_{X}) \to (Y, \taut_{Y})$ una aplicación y $x_0 \in X$. Entonces, los siguientes enunciados son equivalentes:
     \begin{enumerate}
        \item[(1)] $f$ es continua en $x_0$
        \item[(2)] Para cada conjunto infinito $I$, cada filtro sobre $I$ $\filt$ y cada colección de elementos $(x_{\ii})_{\ii \in I} \subseteq X$ tal que $x_{\ii} \underset{\ii \to \filt}{\overset{\taut_{X}}{\to}} x_{0}$, se cumple que $f(x_{\ii}) \underset{\ii \to \filt}{\overset{\taut_{Y}}{\to}} f(x_{0})$.
        \item[(3)] Para cada conjunto infinito $I$, cada ultrafiltro sobre $I$ $\ultf$ y cada colección de elementos $(x_{\ii})_{\ii \in I} \subseteq X$ tal que $x_{\ii} \underset{\ii \to \ultf}{\overset{\taut_{X}}{\to}} x_{0}$, se cumple que $f(x_{\ii}) \underset{\ii \to \ultf}{\overset{\taut_{Y}}{\to}} f(x_{0})$.   
     \end{enumerate}
    \end{prop*}

    \begin{dem}
    
        \borde{(1) \implies (2)} Sea $f$ continua en $x_0$. Tomemos un conjunto infinito $I$, un filtro sobre $I$ $\filt$ y una colección de elementos $(x_{\ii})_{\ii \in I} \subseteq X$ tal que $x_{\ii} \underset{\ii \to \filt}{\overset{\taut_{X}}{\to}} x_{0}$. Como $f$ es continua en $x_{0}$, sabemos que se cumple que:
        \[
        \forall W \in \mathcal{N}_{f(x_0)} \ \exists V \in \mathcal{N}_{x_0} : f(V) \subseteq W
        \]

        \medskip
        Por otro lado, como $x_{\ii} \underset{\ii \to \filt}{\overset{\taut_{X}}{\to}} x_{0}$, se cumple que:
        \[
        I_{V} = \{\ii \in I \colon x_{\ii} \in V\} \in \filt
        \]
        Entonces, juntando ambas deducciones, tenemos que:
        \[
        \forall \ii \in I_{V} \ f(x_{\ii}) \in W \implies \{\ii \in I_{V} : f(x_{\ii} \in W) \in W\} \supset I_{V} \in \filt
        \]
        y por las propiedades de filtro, se sigue que $\{\ii \in I \colon f(x_{\ii}) \in W\} \in \filt$. Es decir, que $f(x_{\ii}) \underset{\ii \to \filt}{\overset{\taut_{Y}}{\to}} f(x_{0})$.

        \borde{(2) \implies (3)} Esta implicación es inmediata, ya que todo ultrafiltro es un filtro.

        \borde{(3) \implies (1)} Supongamos que $f$ no es continua en $x_0$., es decir, que se verifica lo siguiente:
        \[
        \exists W \in \mathcal{N}_{f(x_0)} \ \forall V \in \mathcal{N}_{x_0} : f(V) \not \subseteq W
        \]

        Por lo tanto, para cada $V \in \mathcal{N}_{x_0}$ existe un elemento $x_{V} \in V$ tal que $f(x_{V}) \not \in W$, o expresado de otro modo:
        \[
        \{V \in \mathcal{N}_{x_0} :f(x_{V}) \in W\} = \emptyset
        \]

        \medskip
        Consideramos el conjunto dirigido $(\mathcal{N}_{x_0}, \preceq)$ y el filtro de orden sobre $\mathcal{N}_{x_0}$, $\filt$, el cuál extendemos a un ultrafiltro $\ultf$ que contiene a $\filt$. Por las observaciones del ejemplo previo, se tiene que $x_{V} \underset{V \to \ultf}{\overset{\taut_{X}}{\to}} x_{0}$, pero por construcción, $f(x_{V}) \underset{V \to \ultf}{\overset{\taut_{Y}}{\to}} f(x_{0})$, ya que $\{V \in \mathcal{N}_{x_0} : f(x_{V}) \in W\} = \emptyset \not \in \ultf$
    \end{dem}

    \begin{lem*}{}
    Sean $\ultf$ un ultrafiltro sobre un conjunto infinito $I$, $\{I_{\kk}\}_{\kk = 1}^{\nn}$ una colección finita de subconjuntos de $I$ tal que $\forall \kk \in \{1, \cdots, \nn\} \ I_{\kk} \not \in \ultf$. Entonces:
    \[
    \bigcup_{\kk=1}^{\nn} I_{\kk} \not \in \ultf
    \]
    \end{lem*}

    \begin{dem}
    Distinguiremos dos casos:
    
    \medskip

     Suponemos por reducción al absurdo que $\bigcup_{\kk=1}^{\nn} I_{\kk} \in \ultf$. Por hipótsesis, se verifica lo siguiente:
     \begin{itemize}
        \item[(a)] Por ser $\ultf$ un ultrafiltro, aplicando la caracterización de ultrafiltros, se tiene que $\forall \kk \in \{1, \cdots, \nn\} : I \setminus I_{\kk} \in \ultf$.
        \item[(b)] Por ser $\ultf$ un filtro, es cerrado bajo intersecciones finitas.
     \end{itemize}

     \smallskip
     A partir de estas dos observaciones, tenemos que:
     \[
     \begin{aligned}
        I \setminus \bigcup_{\kk=1}^{\nn} I_{\kk} = \bigcap_{\kk=1}^{\nn} \ (I \setminus I_{\kk})  \in \ultf \textcolor{red}{\#}
     \end{aligned}
     \]
    \end{dem}

    \begin{thm*}{Caracterización de la compacidad en términos de ultrafiltros}
        Sea $(X, \taut)$ un espacio topológico. Entonces, los siguientes enunciados son equivalentes:
        \begin{enumerate}
            \item[(1)] $(X, \taut)$ es compacto
            \item[(2)] Para cada conjunto infinito $I$, cada ultrafiltro sobre $I$ $\ultf$ y cada colección de elementos $(x_{\ii})_{\ii \in I} \subseteq X$, existe un elemento $x_{0} \in X$ tal que $x_{\ii} \underset{\ii \to \ultf}{\overset{\taut}{\to}} x_{0}$.
        \end{enumerate}
    \end{thm*}

    \begin{dem}
    \borde{(1) \implies (2)} Supongamos que $(X, \taut)$ es compacto y que por reducción al absurdo, existen un conjunto de índices $I$, un ultrafiltro sobre $I$, $\ultf$ y una colección de elementos $(x_{\ii})_{\ii \in I} \subseteq X$ tales que no existe ningún $x_0 \in X$ verificando $x_{\ii} \underset{\ii \to \ultf}{\overset{\taut}{\not \to}} x_{0}$.

    \smallskip
    Entonces, para cada $x \in X$, existe un entorno abierto $V_{x} \in \mathcal{N}_{x}$ cumpliendo que:
    \[
    I_{x} \coloneq \{\ii \in I \colon x_{\ii} \in V_{x}\} \not \in \ultf
    \]

    \smallskip
    Por ser $(X, \taut)$ compacto y darse que $X = \bigcup_{x \in X} V_x$, existe un subconjunto finito $\{z_1, z_2, \ldots, z_{\nn}\} \subseteq X$ tal que $X = \bigcup_{\ii=1}^{\nn} V_{z_{\ii}}$.
    
    \smallskip
    De este modo, se tiene lo siguiente:
    \begin{enumerate}
        \item[(a)] Por construcción, $I = \bigcup_{\ii = 1}^{\nn} I_{z_{\ii}}$, ya que dado cualquier índice $\ii \in I$, por ser $\bigcup_{\jj =1}^{\nn} V_{z_{\jj}}$ recubrimiento de $X$, debe esxistir un $\kk \in \{1, \cdots, \nn\}$ tal que $x_{\ii} \in V_{z_{\kk}}$.
        Por lo tanto, $\ii \in I_{z_{\kk}} \subseteq \bigcup_{\jj = 1}^{\nn} I_{z_{\jj}}$ .
        \item[(b)] Por ser $\ultf$ un filtro, se tiene que $I \in \ultf$
        \item[(c)] Por ser $\ultf$ un ultrafiltro y por la construcción de los conjuntos $I_{z_{\ii}}$, se tiene que $I_{z_{\ii}} \not \in \ultf$ para cada $\ii \in \{1, \cdots, \nn\}$. Por el lema previo, se sigue que $\bigcup_{\ii = 1}^{\nn} I_{z_{\ii}} \not \in \ultf$. 
    \end{enumerate}
    
    \medskip
    De los puntos (a), (b) y (c) se deduce una contradicción \textcolor{red}{\#}


    \medskip
    \borde{(2) \implies (1)} Asumiendo la existencia de un límite bajo el ultrafiltro $\ultf$, consideramos una familia de cerrados $\{K_{\ii} \colon \ii \in I\} \in \mathcal{P}(X)$ con la PIF. Por la caracterización de la compacidad en términos de cerrados, basta ver que se cumple:
     \[
     \bigcap_{\ii \in I} K_{\ii} \neq \emptyset
     \]

     \medskip
     Sea $\mathcal{D} = \{D \subset I : D \ \text{es finito}\}$ dotado del orden inclusión de conjuntos $(\mathcal{D}, \subseteq)$. Evidentemente, este conjunto es dirigido, pues dados $D, J \in \mathcal{D}$, resulta claro que $D \cup J \in \mathcal D$ junto con $D \subseteq D \cup J$ y $J \subseteq D \cup J$. Por lo tanto, podemos considerar el filtro de orden sobre $\mathcal{D}$, $\filt$, y extenderlo a un ultrafiltro $\ultf$ que contenga a $\filt$.

     \smallskip
     Para cada $D \in \mathcal{D}$, por tener la PIF, se tiene que $\bigcap_{\ii \in D} K_{\ii} \neq \emptyset$, de modo que
     \[
    \exists \ x_{D} \in \bigcap_{\ii \in D} K_{\ii}
     \]

    \smallskip
    Como asumimos que $\exists \underset{D \to \ultf}{\lim x_{D}} = x_{0} \in X$, trataremos de probar que $x_{0} \in \bigcap_{\ii \in I} K_{\ii}$.
    
    \smallskip
    Fijamos un $\ii \in I$ cualquiera y tomamos un entorno $V \in \mathcal{N}_{x_0}$. Probemos que $V \cap K_{\ii} \neq \emptyset$.
    \[
    \left.
    \begin{aligned}
        x_{D} \underset{D \to \ultf}{\overset{\taut}{\to}} x_{0} &\implies \mathcal{A}_{V} \coloneq \{D \in \mathcal{D} : x_{D} \in V\} \in \ultf\\
         \mathcal{D}(\{\ii\}) \in \ultf
    \end{aligned}
    \right\} \implies \mathcal{A}_{V} \cap \mathcal{D}(\{\ii\}) \underset{\ultf \text{filtro}}{\neq} \emptyset
    \]

    Por lo tanto, existe un conjunto finito $J \in \mathcal{A}_{V} \cap \mathcal{D}(\{\ii\})$, que verifica
    \[
    \left.
    \begin{aligned}   
    J \in \mathcal{A}_{V} & \implies x_{J} \in V\\
    J \in \mathcal{D}(\{\ii\}) & \implies J \in \mathcal D \text{ y } \{\ii\} \subset J \implies \ii \in J \underset{x_J \in \cap_{\ii \in J} K_{\ii}}{\implies} x_J \in K_{\ii}
    \end{aligned}
    \right\} \implies x_J \in V \cap K_{\ii} \neq \emptyset
    \]

    \medskip
    Por la definición de clausura, se sigue que $\forall \ii \in I \ : x_{0} \in \overline{K_{\ii}} \underset{K_{\ii} \ \text{cerrados}}{=} K_{\ii}$, es decir:
    \[
    x_{0} \in \bigcap_{\ii \in I} K_{\ii} \neq \emptyset
    \]
    \end{dem}


    \begin{prop*}{}
        Sean $\{(X_{\alpa}, \taut_{\alpa}) \ \colon \alpa \in \Delt\}$ una familia de espacios topológicos, $I$ un conjunto de índices, $\ultf$ un ultrafiltro sobre $I$,una familia $(\mathbb{X}_{\ii})_{\ii \in I} \subseteq \prod_{\alpa \in \Delt} X_{\alpa}$ y $\mathbb{X}_{0} \in \prod_{\alpa \in \Delt} X_{\alpa}$. Entonces, los siguientes resultados son equivalentes:
        \begin{enumerate}
            \item[(1)] $\mathbb{X}_{\ii} \underset{\ii \to \ultf}{\overset{\prod_{\alpa \in \Delt} \taut_{\alpa}}{\to}} \mathbb{X}_{0}$ 
            \item[(2)] Para cada $\alpa \in \Delt$, se cumple que $x_{\ii}^{\alpa} \underset{\ii \to \ultf}{\overset{\taut_{\alpa}}{\to}} x_{0}^{\alpa}$, donde $x_{\ii}^{\alpa} = p_{\alpa}(\mathbb{X}_{\ii})$ y $x_{0}^{\alpa} = p_{\alpa}(\mathbb{X}_{0})$.
        \end{enumerate}     
    \end{prop*}

    \begin{dem}
        \borde{(1) \implies (2)} Consideremos un $\berta \in \Delt$ cualquiera pero fijo. Definimos la aplicación proyección canónica, trivialmente continua en $\prod_{\alpa \in \Delt} \taut_{\alpa}$:

        \bigskip
        \[
        \begin{aligned}
            p_{\berta} \colon  & \prod_{\alpa \in \Delt} X_{\alpa} \to X_{\berta}\\
            & (x^{\alpa})_{\alpa \in \Delt} \mapsto x^{\berta}
        \end{aligned}
        \]
        
        \medskip
        Por la caracterización de la continuidad en términos de filtros junto con la hipótesis $\mathbb{X}_{\ii} \underset{\ii \to \ultf}{\overset{\prod_{\alpa \in \Delt} \taut_{\alpa}}{\to}} \mathbb{X}_{0}$, se sigue que:

        \[
        p_{\berta}(\mathbb{X}_{\ii}) \underset{\ii \to \ultf}{\overset{\taut_{\berta}}{\to}} p_{\berta}(\mathbb{X}_{0}) \iff x_{\ii}^{\berta} \underset{\ii \to \ultf}{\overset{\taut_{\berta}}{\to}} x_{0}^{\berta}
        \]
    \end{dem}


    \section{Ejercicios}
     
    \begin{ejer*}{1}
        Sea $\filt$ el filtro de Fréchet sobre $\NN$ y sea $\filt_{\nn}$ el filtro principal sobre $\NN$ asociado al número natural $\nn$
        \smallskip
        \begin{enumerate}
            \item[(a)] ¿Qué filtro es $\gilt = \bigcap_{n \in \NN} \filt_{n}$?
            \item[(b)] ¿Es $\gilt \cap \filt $ un filtro sobre $\NN$?
        \end{enumerate}
    \end{ejer*}

    \begin{resol}
        Veamos cada uno de los apartados:

        \medskip
        \borde{(a)} Por definición, sabemos que el filtro principal para cada número natural viene dado pir $\filt_{\nn} = \{A \subseteq \NN \colon \nn \in \NN\}$. De esta forma. al considerar la intersección, se tiene que:
        \[
        \begin{aligned}
            A \in \bigcap_{\nn \in \NN} \filt_{\nn} &\iff A \in \filt_{\nn} \ \forall \nn \in \NN\\
            &\iff \forall \nn \in \NN \colon \nn \in A \iff A = \NN
        \end{aligned}
        \]

        Luego, hemos deducido que $\gilt$ se trata del filtro trivial sobre $\NN$

        \medskip
        \borde{(b)} Por la deducción previa, tenemos que $\gilt = \{\NN\}$. Por lo tanto, $\gilt \cap \filt = \{\NN\} \cap \filt = \{\NN\}$, que sigue siendo un filtro sobre $\NN$.
    \end{resol}

    \begin{ejer*}{2}
        Sea $(x_{\nn})_{\nn \in \NN}$ una sucesión convergente a $x_0$ en un espacio topológico.

        \smallskip
        Demuestra que toda reordenación suya $(x_{\nn})_{\nn \in \NN}$ converge a $x_0$. (Entenderemos por reordenación a una sucesión $(x_{\varphi(\nn)})_{\nn \in \NN}$ donde $\varphi : \NN \to \NN$ es una función biyectiva.)

    \end{ejer*}

     \begin{resol}
    Consideramos un espacio topológico $(X, \taut)$ y una sucesión $(x_{\nn})_{\nn \in \NN} \subseteq X$ convergente a $x_0 \in X$. Por definición de convergencia, 
     \end{resol}


    \begin{ejer*}{6 - \ Teorema de Ramsey}
        Dado un conjunto $J$, denotamos $\mathcal{P}(J)$ al conjunto de todos los subconjuntos de $J$ y $\mathcal{P}^{\kk}(J)$ al conjunto de todos los subconjuntos de cardinal $\kk$ de $J$.

        \medskip
        Sean $A$ y $B$ dos subconjuntos disjuntos de $\mathcal{P}^{2}(\NN)$ tales que $\mathcal{P}^{2}(\NN) = A \cup B$.
        
        \smallskip
        Demuestra que existe un subconjunto infinito $I \subseteq \NN$ tal que o bien $\mathcal{P}^{2}(I) \subseteq A$ o bien $\mathcal{P}^{2}(I) \subseteq B$.
    \end{ejer*}

    \begin{resol}
        \borde{Dem \ 1} Tomemos un número natural $\nn_{0} \in \NN$. Sabemos que existen infinitas posibles duplas a formar con $\nn_{0}$, es decir, existen infinitos elementos de la forma $\{\nn_{0}, \nn\}$ con $\nn \in \NN \setminus \{\nn_{0}\}$. Por lo tanto, como $\mathcal{P}^{2}(\NN) = A \cup B$, debe existir un subconjunto infinito $\mathcal{N}_{1} \subseteq\NN $ tal que o bien $\{\nn_{0}, \nn\} \in A$ para todo $\nn \in \mathcal{N}_{1}$ o bien $\{\nn_{0}, \nn\} \in B$ para todo $\nn \in \mathcal{N}_{1}$.

        \medskip
        Independientemente de cuál de los dos casos se cumpla, podemos tomar ahora un elemento $\nn_1 \in \mathcal{N}_{1}$. Por ser un conjunto infinito, existen infinitos elementos de la forma $\{\nn_{1}, \nn\}$ con $\nn \in \mathcal{N}_{1} \setminus \{\nn_{1}\}$. Por lo tanto, como $\mathcal{P}^{2}(\NN) = A \cup B$, debe existir un subconjunto infinito $\mathcal{N}_{2} \subseteq\mathcal{N}_{1} $ tal que o bien $\{\nn_{1}, \nn\} \in A$ para todo $\nn \in \mathcal{N}_{2}$ o bien $\{\nn_{1}, \nn\} \in B$ para todo $\nn \in \mathcal{N}_{2}$.

        \medskip
        Reiterando este proceso indefinidamente, obtenemos una colección de naturales $(\nn_{\kk})_{\kk \in \NN} \subseteq \NN$ y una colección de subconjuntos infinitos $(\mathcal{N}_{\kk})_{\kk \in \NN} \subseteq \mathcal{P}(\NN)$. 
        Ahora bien, podemos considerar la subsucesión $(\nn_{\kk}^{A})_{\kk \in \NN}$ a partir de los conjuntos $(\mathcal{N}_{\kk}^{A})_{\kk \in \NN}$  cuyos elementos tienen infinitos pares en $A$. De esta forma, si tomamos el conjunto infinito $I =(\nn_{\kk}^{A})_{\kk \in \NN} $, es claro que $\mathcal{P}^{2}(I) \subseteq A$. De forma análoga, podemos replicar el argumento para los conjuntos $(\mathcal{N}_{\kk}^{B})_{\kk \in \NN}$.

        \bigskip
        \borde{Dem \ 2} Consideramos un ultrafiltro $\ultf$ sobre $\NN$ que contenga al filtro de Fréchet sobre $\NN$, $\filt$. Para cada $\nn \in \NN$, por hipótesis, podemos definir los conjuntos:
        \[
        \begin{aligned}
            S_{A}(\nn) &= \{\mm \in \NN : \{\nn, \mm\} \in A\}\\
            S_{B}(\nn) &= \{\mm \in \NN : \{\nn, \mm\} \in B\}
        \end{aligned}
        \]

        \medskip
        Por ser $\ultf$ un ultrafiltro que contiene al filtro de Fréchet, para cada $\nn \in \NN$ se tiene que $S_{A}(\nn) \cup S_{B}(\nn) = \NN \setminus \{\nn\} \in \filt \subseteq \ultf$.
        
        \smallskip
        Por el recíproco de un lema sobre ultrafiltros, sabemos que al menos uno de los dos conjuntos $S_{A}(\nn)$ o $S_{B}(\nn)$ pertenece a $\ultf$. Esto nos permite definir los conjuntos:
        \[
        \begin{aligned}
            \ultf_{A} &= \{\nn \in \NN : S_{A}(\nn) \in \ultf\}\\
            \ultf_{B} &= \{\nn \in \NN : S_{B}(\nn) \in \ultf\}
        \end{aligned}
        \]

        \medskip
        Por un razonamiento análogo al anterior, como $\ultf_{A} \cup \ultf_{B} = \NN$, al menos uno de los conjuntos $\ultf_{A}$ o $\ultf_{B}$ pertenece a $\ultf$.
        
        \medskip
        Supongamos sin perdida de generalidad que $\ultf_{A} \in \ultf$ (el caso de $\ultf_{B}$ es análogo). Entonces, si tomamos un elemento $\nn_{0} \in \ultf_{A}$, se tiene que $S_{A}(\nn_{0}) \in \ultf$.

        \smallskip
        Seguidamente, como $\ultf_{A} \in \ultf$ y $S_{A}(\nn_{0} \in \ultf)$, se tiene que $\ultf_{A} \cap S_{A}(\nn_0) \neq \emptyset$ y de hecho, podemos asegurar que se trata de un conjunto infinito, pues $\ultf$ extiende al filtro de Fréchet $\filt$ que solo contien sunconjuntos numerables.
        De esta forma, podemos tomar un elemento $\nn_{1} \in \ultf_{A} \cap S_{A}(\nn_0)$ y repetir el proceso de forma iterativa, obteniendo una colección de elementos $(\nn_{\kk})_{\kk \in \NN} \subseteq \NN$ tal que $\mathcal{P}^{2}(I) \subseteq A$ para $I = (\nn_{\kk})_{\kk \in \NN}$.
    \end{resol}


    \begin{ejer*}{7}
        Sean $(X, \taut_{X})$ e $(Y,\taut_{Y})$ dos espacios espacios topológicos con $(X, \taut_{X})$ compacto y $f : (X,\taut_{X}) \to (Y,\taut_{Y})$ una aplicación contínua entre ellos.

        \smallskip
        Demuestra, empleando filtros y ultrafiltros, que $f(X)$ es compacto en $(Y,\taut_{Y})$.

    \end{ejer*}

    \begin{resol}
    Para realizar este ejercicio, emplearemos las caracterizaciones de la continuidad y compacidad via ultrafiltros. 
    
    \medskip
    Así, tomamos un conjunto de índices $I$, un ultrafiltro sobre $I$ $\ultf$ y una colección de elementos $(y_{\ii})_{\ii \in I} \subseteq f(X)$. Queremos probar que existe un elemento $y_{0} \in f(X)$ verificando que:
    \[
    y_{\ii} \underset{\ii \to \ultf}{\overset{\taut_{Y}}{\to}} y_{0}
    \]

    \smallskip
    Sin embargo, por construcción, para cada $\ii \in I$, existe un elemento $x_{\ii} \in X$ verficando que $y_{\ii} = f(x_{\ii})$. A tenor de esto, aplicamos sendas caracterizaciones:
    \begin{enumerate}
        \item[(1)] Por la caracterización de la compacidad en términos de ultrafiltros, como $(X,\taut_{X})$ es compacto, sabemos que para la familia $(x_{\ii})_{\ii \in I} \subseteq X$ existe un elemento $x_{0} \in X$ tal que $x_{\ii} \underset{\ii \to \ultf}{\overset{\taut_{X}}{\to}} x_{0}$
        \item[(2)] Por la caracterización de la continuidad en términos de ultrafiltros, como $f$ es continua, sabemos que $f(x_{\ii}) \underset{\ii \to \ultf}{\overset{\taut_{Y}}{\to}} f(x_{0})$
        \item[] 
    \end{enumerate}

    Es decir, si definimos $y_{0} = f(x_{0})$, entonces $y_{\ii} \underset{\ii \to \ultf}{\overset{\taut_{Y}}{\to}} y_{0}$. Aplicando nuevamente la caracterización de la compacidad en términos de ultrafiltros, tenemos que $f(X)$ es compacto en $(Y,\taut_{Y})$. 
    \end{resol}

    \begin{ejer*}{8}
       Sean $(X, \taut_{X})$ e $(Y,\taut_{Y})$ dos espacios espacios topológicos con $(X, \taut_{X})$ compacto y $f : (X,\taut_{X}) \to (Y,\taut_{Y})$ una aplicación biyectiva y contínua entre ellos.

        \smallskip
        Demuestra, empleando filtros y ultrafiltros, que $f^{-1}$ es una aplicación contínua.
        
    \end{ejer*}

    \begin{resol}
        Para realizar este ejercicio, emplearemos las caracterizaciones de la continuidad y la compacidad en términos de ultrafiltros, además del ejercicio previo.

        \medskip
        Tomamos un conjunto de índices $I$, $\ultf$ un ultrafiltro sobre $I$ y una colección de elementos $(y_{\ii \in I})_{\ii \in I} \subseteq Y$ verificando que $y_{\ii \in I} \underset{\ii \to \ultf}{\overset{\taut_{Y}}{\to}} y_{0}$ para cierto $y_0 \in Y$. Queremos probar que:
        \[
        f^{-1}(y_{\ii}) \underset{\ii \to \ultf}{\overset{\taut_{X}}{\to}} f^{-1}(y_{0})
        \]

        \medskip
        Por ser $f$ es biyectiva, sabemos que existe una colección de elementos  $(x_{\ii})_{\ii \in I} \subseteq X$ tal que $f(x_{\ii}) = y_{\ii}$. Por la caracterización de la compacidad en términos de ultrafiltros, debe existir un elemento $x_0 \in X$ tal que $x_{\ii} \underset{\ii \to \ultf}{\overset{\taut_{X}}{\to}} x_{0}$.

        \smallskip
        Por otro lado, como $f$ es continua, se tiene que por la caracterización de la continuidad en términos de ultrafiltros, $f(x_{\ii}) \underset{\ii \to \ultf}{\overset{\taut_{Y}}{\to}} f(x_{0})$.

        \smallskip
        Ahora bien, como estamos trabajando en espacios Hausdorff, por construcción de las sucesiones deben coincidir: $f(x_0) = y_0$. De esta forma, recordando la biyectividad, se tiene wue $f^{-1}(y_0) = x_0$ y $(x_{\ii})_{\ii \in I} = (f^{-1}(y_{\ii}))_{\ii \in I}$. Luego, hemos probado lo que buscábamos
    \end{resol}

    \begin{ejer*}{Unicidad de límites en espacios Hausdorff}
        Sea $(X, \taut)$ un espacio Hausdorff y sea $\filt$ un filtro sobre un conjunto de índices $I$.
        
        \smallskip
        Demuestra que si existe un elemento $x_0 \in X$ tal que $x_{\ii} \underset{\ii \to \filt}{\overset{\taut}{\to}} x_{0}$, entonces $x_0$ es único.
    \end{ejer*}

    \begin{resol}
        Supongamos por reducción al absurdo que existen dos elementos distintos $x_0, x_1$ que son límite de la colección $(x_{\ii})_{\ii \in I}$ bajo el filtro $\filt$. Entonces, se cumple lo siguiente:
        \begin{enumerate}
            \item[(1)] Por hipótesis, sabemos que para cada entorno abierto $V_0 \in \mathcal{N}_{x_0}$, se tiene que $I_{V_0} \coloneq \{\ii \in I \colon x_{\ii} \in V_0\} \in \filt$.
            \item[(2)] De forma análoga, para cada entorno abierto $V_1 \in \mathcal{N}_{x_1}$, se tiene que $I_{V_1} \coloneq \{\ii \in I \colon x_{\ii} \in V_1\} \in \filt$.
            \item[(3)] Por hipótesis, como $(X, \taut)$  es Hausdorff, existen entornos abiertos $V_{0}^{*} \in \mathcal{N}_{x_0}$ y $V_{1}^{*} \in \mathcal{N}_{x_1}$ tales que $V_{0}^{*} \cap V_{1}^{*} = \emptyset$.
        \end{enumerate}

        A partir de estas observaciones, tenemos que:
        \[
        \left.
        \begin{aligned}
            I_{V_{0}^{*}} \in \filt \\
            I_{V_{1}^{*}} \in \filt
        \end{aligned}
        \right\} \implies I_{V_{0}^{*}} \cap I_{V_{1}^{*}} \in \filt
        \]

        Sin embargo, por como se han tomado los entornos $V_{0}^{*}$ y $V_{1}^{*}$, se tiene que $I_{V_{0}^{*}} \cap I_{V_{1}^{*}} = \emptyset$, \textcolor{red}{\# lo cual contradice la definición de filtro}. 
        
        \medskip
        Por lo tanto, hemos probado que el límite ha de ser único..
    \end{resol}

    \begin{ejer*}{Importancia }
        Sea $(X, \tau)$ un espacio topológico \textbf{\textit{primero numerable}} y sea $\mathcal{F}$ un filtro sobre un conjunto de índices $I$. Entonces, los siguientes enunciados son equivalentes:
        \begin{enumerate}
    \end{ejer*}


% ======================================================================= %
%                              ZONA DE ANEXOS                             %
% ======================================================================= %
% ======================================================================= %
%                              BIBLIOGRAFÍA                               %
% ======================================================================= %
\cleardoublepage
\printbibliography[heading=bibintoc]

\newpage
\printindex
\end{document}